\documentclass[11pt]{article}
\usepackage[T1]{fontenc}
\usepackage{lmodern}
\usepackage[margin=1in]{geometry}
\usepackage{amsmath,amssymb,amsthm,mathtools,mathrsfs}
\usepackage{microtype}
\usepackage{enumitem}
\usepackage[colorlinks=true,linkcolor=blue,citecolor=blue,urlcolor=blue]{hyperref}
\hypersetup{pdftitle={Abelian covers, Gale correspondences, and the Hodge conjecture for powers},pdfauthor={OpenAI}}
\newcommand{\Q}{\mathbb Q}
\newcommand{\C}{\mathbb C}
\newcommand{\Z}{\mathbb Z}
\newcommand{\Qbar}{\overline{\mathbb Q}}
\newcommand{\PP}{\mathbb P}
\newcommand{\A}{\mathbb A}
\newcommand{\Gm}{\mathbb G_m}
\newcommand{\CH}{\operatorname{CH}}
\newcommand{\Gr}{\operatorname{Gr}}
\newcommand{\Jac}{\operatorname{Jac}}
\newcommand{\Conf}{\operatorname{Conf}}
\newcommand{\pr}{\operatorname{pr}}
\newcommand{\id}{\operatorname{id}}
\newcommand{\kk}{\overline{\Q}}
\newcommand{\U}{\mathcal U}
\newcommand{\D}{\mathcal D}
\newcommand{\V}{\mathbb V}
\newcommand{\cY}{\mathcal Y}
\theoremstyle{plain}
\newtheorem{theorem}{Theorem}[section]
\newtheorem{proposition}[theorem]{Proposition}
\newtheorem{lemma}[theorem]{Lemma}
\newtheorem{corollary}[theorem]{Corollary}
\theoremstyle{definition}
\newtheorem{definition}[theorem]{Definition}
\newtheorem{companionresults}[theorem]{Companion results}
\theoremstyle{remark}
\newtheorem{remark}[theorem]{Remark}
\numberwithin{equation}{section}
\title{Abelian covers, Gale correspondences,\\
and the Hodge conjecture for powers}
\author{OpenAI}
\date{September 30, 2026}
\begin{document}
\maketitle

\begin{abstract}
We prove the rational Hodge conjecture on every self-power of the
Jacobian at each tensor Hodge-generic point of the full marked variation
of a connected abelian cover of curves.  This holds in every base genus
and for every compatible branching pattern.  For any CM abelian variety,
the conclusion also holds for every self-power of its product with the
Jacobian, on the same Hodge-generic locus.
We also prove the conjecture on every self-power of a very general
member of the full smooth labelled family of diagonal complete
intersections cut out by at most two equations of a common degree, in
every dimension and every degree at least two.
\end{abstract}

\tableofcontents

\section{Introduction}

For a smooth projective complex variety \(X\), the rational Hodge
conjecture asserts that the cycle class map
\[
 \CH^p(X)_\Q\longrightarrow
 H^{2p}(X,\Q)\cap H^{p,p}(X,\C)
\]
is surjective for every \(p\).  Its assertion for all self-powers
\(X^u\) controls classes coupling different copies of \(X\) and is
stronger than the assertion for \(X\) alone.  We prove all-self-power
results for two families.  The first concerns Jacobians of connected
abelian covers of curves.  The second concerns diagonal complete
intersections; their cohomology is supplied by correspondences from
the cover Jacobians and CM abelian varieties.

\subsection{Abelian covers}

Fix a compact connected oriented surface \(B_{\mathrm{top}}\), a finite
labelled set \(D\) of branch points, a point \(*\notin D\), a finite
abelian group \(P\), and an epimorphism
\[
 \phi:\pi_1(B_{\mathrm{top}}\setminus D,*)\longrightarrow P
\]
whose branch meridians have nonidentity image.  The associated cover is
connected; denote its compact covering surface by \(C_{\mathrm{top}}\)
and retain a chosen point over \(*\).  Empty branching and the
trivial group are allowed.  Vary the full marked complex structure of
the pointed base, including all branch positions, and give the cover
the induced complex structure.  Add labelled unbranched auxiliary
points only when needed for stability.  Denote the connected marked
parameter manifold by \(\mathcal T_\phi\), its fibers by \(C_s\), and its marked
rational first cohomology by
\(H_\Q=H^1(C_{\mathrm{top}},\Q)\).

Let \(\mathcal T_\phi^{\mathrm{Hg}}\) consist of the parameters \(s\)
with the following property: for every \(a\geq0\), any rational tensor
in \(H_\Q^{\otimes 2a}\) of type \((a,a)\) at \(s\) has that type
throughout \(\mathcal T_\phi\).  Lemma~\ref{p4:generic-tensors} proves
that this is the complement of a countable union of proper closed
analytic tensor Hodge loci.  A CM abelian variety here means one whose
degree-one Hodge group is a torus; a point is allowed.

\begin{theorem}[Powers of abelian-cover Jacobians]
\label{p4:abelian-covers}
For each fixed connected topological abelian cover as above, every
\(s\in\mathcal T_\phi^{\mathrm{Hg}}\), every complex CM abelian variety
\(M\), and every integer \(N\geq1\), the rational Hodge conjecture holds
in every codimension of
\[
 \bigl(\Jac(C_s)\times M\bigr)^N.
\]
It also holds for \(C_s^a\times M^b\) for all integers \(a,b\geq0\),
where a zeroth power means a point.  The point \(s\) is chosen before
\(M\): the same Hodge-generic locus works for every CM companion.
\end{theorem}

The theorem allows every base genus and every branching pattern
compatible with the fixed connected cover.  Its genericity is taken in
the full marked variation.  Applying it to a constrained algebraic
family requires a comparison with that variation and a check that its
exceptional tensor loci pull back properly.  The second part of this
paper carries out both steps for the curves arising from Gale duality.

\subsection{Diagonal complete intersections}

Fix an integer \(n\geq2\).  For a rank-\(c\) matrix
\(A=(a_{\alpha i})\) with \(m\) columns, let
\[
 X_A=\left\{[x_1:\cdots:x_m]\in\PP^{m-1}:
       \sum_{i=1}^m a_{\alpha i}x_i^n=0
       \quad(1\leq\alpha\leq c)\right\}.
\]
Changing the row basis of \(A\) does not change \(X_A\).  We parametrize
these intersections by the row space in \(\Gr(c,m)\), retaining the
labels of the projective coordinates.  Write
\[
 \D_{c,m}=\{Q\in\Gr(c,m):\text{every maximal coordinate minor
 of \(Q\) is nonzero}\}.
\]
Proposition~\ref{prop:smooth-locus} proves that this is exactly the
smooth locus when \(1\leq c\leq m-1\).  The dimension there is
\(d=m-c-1\).  For \(c=0\) the parameter space is a point and
\(X=\PP^{m-1}\).

\begin{theorem}[Powers of diagonal complete intersections]
\label{thm:main}
Fix \(n\geq2\),
\(d\geq0\), \(c\in\{0,1,2\}\), and put \(m=d+c+1\).
Outside a countable union of proper closed algebraic subsets of
\(\D_{c,m}\), the rational Hodge conjecture holds in every codimension
on \(X_A^u\) for every \(u\geq1\).  For \(c\leq1\), or for \(d=0\),
no exceptional subset is needed.
\end{theorem}

This is the full smooth family of labelled diagonal presentations of
the indicated rank, not the moduli space of the underlying abstract
varieties.  The Gale construction itself applies at every smooth
parameter.  The very-general restriction enters only when its dual
curve is placed in the locus of Theorem~\ref{p4:abelian-covers}.

\subsection{Earlier work}

The determinant construction for curves has direct precedent in
Schoen's phase-kernel and permutation quotient
\cite[Section~1, Lemmas~1.1--1.2]{Schoen}.  His algebraicity theorem
there treats an even-rank setting with a pairable branch tuple
\cite[Theorem~2.0]{Schoen}.  His construction also uses the distinguished
Abel fiber \cite[Section~2, pp.~12--13]{Schoen}.  We isolate the top
exterior character line by this quotient and obtain arbitrary-branch
correspondence supply by dominating the distinguished fiber with
bundles over Fermat hypersurfaces.
The Fermat domination is due to Shioda and Katsura~\cite{SK}.  The supply
uses full cohomology of CM sources; an individual unbalanced determinant
need not itself be a Hodge class.

The character local systems and their sphere signatures belong to the
framework of Deligne and Mostow
\cite[Corollary~2.21 and Section~2.23]{DeligneMostow1986}.
For cyclic covers of the projective line, Rohde proves the full
special-unitary connected monodromy of each mixed non-real character eigenspace
\cite[preprint, Theorem~5.1.1]{Rohde2009}; Moonen records these
projections, including the symplectic order-two case, in his study of
special subvarieties \cite[Proposition~4.7]{Moonen2010}.

The Kodaira--Spencer period differential is classical
\cite[Theorem~1.27]{Griffiths1968PeriodsII}; we give the equivariant
calculation needed also when the covering curve has genus one.  For
mapping-class representations of covers, Landesman, Litt, and Sawin
identify connected monodromy with the derived group of the symplectic
centralizer in high-genus ranges, including base genus at least four
for abelian deck groups \cite[Theorem~1.3]{LandesmanLittSawin2025}.
Our proof gives the required all-genus character projections and
within-cover links directly; the low-genus comparisons do not follow
from that theorem.  The use of globally Hodge tensors to control
monodromy is closely related to Andr\'e's argument
\cite[Section~4, Lemma~4]{Andre1992FixedPart}.

The second geometric construction uses classical Gale association:
linear orthogonality modulo diagonal rescaling for labelled projective
configurations.  It is treated systematically by
Dolgachev--Ortland~\cite{DO} and Eisenbud--Popescu~\cite{EP}; the latter
identifies association by complementary linear series
\cite[preprint, Corollary~3.2]{EP}.
Earlier work on complete intersections of Fermat type includes
Aoki~\cite{AokiCI} and Terasoma~\cite{Terasoma}.  We give the
character-level algebraic correspondence argument used here in full,
including its support conditions, zero-dimensional endpoints, and Tate
shifts.  Its evaluation quotient uses the same graded
alternating-character mechanism as Schoen's quotient.

The dual curve in codimension two is a generalized Fermat curve.
Its diagonal model, abelian deck group, and normalized branch
parameters occur in Gonz\'alez-Diez--Hidalgo--Leyton~\cite{GDL} and
Hidalgo--Kontogeorgis--Leyton-\'Alvarez--Paramantzoglou~\cite{HKLP}.
In our notation, the moduli discussion in
\cite[author proof, Section~4.2]{GDL} assumes \((m-2)(n-1)>2\).
The Hodge application requires the relative family with a fixed based
deck datum, not only its fiberwise model.  Section~\ref{sec:family}
supplies the finite coordinate-root base change, the based lift, and
the pathwise comparison with every marked fiber, including the
genus-zero and elliptic boundary cases.

The Fermat factors in the Gale correspondence are essential.  Classical
Hodge results for Fermat varieties include Aoki's proof of the rational
Hodge conjecture for individual Fermat varieties of prime-power degree
\cite[Corollary~2-3]{AokiFermat}.  Our all-powers deduction instead uses
full-cohomology source maps from CM abelian varieties, combined with
Theorem~\ref{p4:abelian-covers} for mixed powers with one cover Jacobian.

\subsection{The proof in outline}

The common Hodge input is the rational Hodge theorem for every complex
CM abelian variety \cite[Theorem~1.1]{OpenAICM}.  The shared tensor tools
from the companion paper on abelian powers reduce balanced degree-one
tensors to pairings and top exterior lines.  Pairings are algebraic
divisor classes.  A determinant line is \emph{CM-supplied} when it lies
in the span of images of full cohomology groups of CM abelian varieties
under rational algebraic correspondences, with the required Tate
shifts.  A rational Hodge class in that span has a rational Hodge lift
to the sources and is therefore algebraic.  Section~\ref{sec:shared}
records exactly the definitions and companion statements used here.

Section~\ref{p4:sec:determinants} supplies each character determinant
for an abelian cover.  A nontrivial character passes to its faithful
cyclic quotient; the invariant character uses the polarization.  On a
finite quotient of a curve power, the relevant character cohomology is
the top alternating line.  An Abel map concentrates it on one fiber,
which is dominated by varieties whose full cohomology is CM-supplied.
Deligne's weight theorem
detects the line on these sources, and a correspondence returns it to
the curve power.  The same section supplies full Fermat cohomology
from CM abelian varieties.

Section~\ref{p4:sec:covers} supplies the other input to the tensor
criterion: the derived Hodge action at the marked Hodge-generic points.
Disk and handle twists force standard character projections.
Comparisons of the full mapping-class actions identify the spaces
supported by the same simple Hodge ideal, compatibly with Hodge types
at every scalar conjugate.  These identifications and the supplied
determinants give saturated data, so the tensor criterion remains
applicable after adjoining any CM abelian variety.

For the diagonal theorem, assume \(c\in\{1,2\}\); projective space is
handled directly.  For a subspace \(R\subset\C^r\), write
\[
 X_R=\{[x_i]\in\PP^{r-1}:(x_1^n,\ldots,x_r^n)\in R\}.
\]
Put \(P=\ker A\), so \(X_P=X_A\).  For a character of the phase
group \((\mu_n)^m/\mu_n\), its support is the set of coordinates
with nonzero exponent.  Section~\ref{sec:roots} shows that projection
to this support preserves the corresponding nontrivial cohomology.

Write \(F_n^a=\{z_1^n+\cdots+z_{a+2}^n=0\}\subset\PP^{a+1}\).
For the ordinary bilinear pairing, let \(Q=P^\perp\).  Coordinate
multiplication defines
\[
 X_P\times X_Q\longrightarrow F_n^{m-2},\qquad
 ([x_i],[y_i])\longmapsto[x_i y_i].
\]
Section~\ref{sec:gale} proves that a full-support Fermat character line
pulls back to a perfect tensor between the two character spaces.
At planes obtained by evaluating polynomials at distinct points, it is
the exterior coproduct on one curve's first cohomology; constancy of
rank then gives perfectness at every
smooth parameter.  For a general support \(S\), contraction gives a
forward rational correspondence from
\[
 F_n^{|S|-2}\times X_{(\pr_S P)^\perp}
\]
whose image contains the corresponding part of \(H^d(X_P)\).
These are maps on whole rational cohomology degrees.  If \(c=2\),
the only possible positive-dimensional dual source is the connected curve
\(C=X_{P^\perp}\); all other dual sources, and every dual source for
\(c=1\), are finite.  Each source retains its Fermat factor.

Section~\ref{sec:family} identifies the coefficient quotient with the
full ordered branch-configuration space.  Finite coordinate-root
choices identify the algebraic families; configuration paths with
transported based lifts reach every marked cover fiber.  Tensor
Hodge loci then give the algebraic very-general set on which the
dual curve satisfies Theorem~\ref{p4:abelian-covers}.

Finally, the Fermat correspondences and the Abel map replace all source
factors by products \(\Jac(C)^v\times M\) with \(M\) CM; in the
one-equation case the Jacobian factor is absent.  Section~\ref{sec:assembly}
checks the degree shifts and applies the shared correspondence-transfer
lemma.  The cover locus is independent of \(M\), so the same parameter
works for every self-power.

\section{Shared algebraic and tensor tools}
\label{sec:shared}

This section fixes the notation needed by the determinant construction
and records the precise companion results used in its Hodge-theoretic
assembly.  Their proofs belong to the companion paper
\cite[Section~2]{OpenAIAbelianPowers}.
No result on split Weil classes is used in this paper.

All varieties are complex, cohomology is singular cohomology, and
algebraic correspondences have rational coefficients unless stated
otherwise.  Put \(k=\Qbar\subset\C\).  A scalar automorphism of \(k/\Q\)
acts on coefficients in the cohomology of the fixed complex variety;
it does not conjugate the variety.  Our Tate convention is
\(V(r)^{p,q}=V^{p+r,q+r}\), so \(V(r)\) has weight
\(\operatorname{wt}(V)-2r\).  A class is algebraic over \(k\) if it is
in the \(k\)-linear span of rational algebraic cycle classes.
For a rational vector space \(V\), write \(V_k=V\otimes_\Q k\).

For a smooth projective variety \(X\), define
\[
 \mathcal C^d(X)=
 \sum_{B,j,\Gamma}\operatorname{im}\!\left(
 \Gamma_*:H^{d-2j}(B,\Q)\longrightarrow H^d(X,\Q)\right),
 \qquad
 \Gamma\in\CH^{\dim B+j}(B\times X)_\Q ,
\]
where \(B\) ranges over CM abelian varieties, including a point, and
\(j\) ranges over integers for which the displayed groups and cycles exist.
The displayed map is a Hodge map from \(H^{d-2j}(B,\Q)(-j)\).
A subspace of \(H^d(X,k)\) is \emph{CM-supplied} if it lies in
\(\mathcal C^d(X)\otimes_\Q k\).  Thus the definition uses the full
cohomology of the sources.  The CM Hodge theorem and the equivalence of
this toric CM convention with its convention are
\cite[Theorem~1.1 and Lemma~2.3]{OpenAICM}.

Let \(X\) be an abelian variety or a smooth connected projective curve.
A subspace \(W\subset H^1(X,k)\) is \emph{admissible} if it is the image
of a \(k\)-linear idempotent whose every scalar conjugate preserves the
Hodge decomposition after extension to \(\C\).  A map between such
spaces, or between finite external direct sums of them, is admissible
if it preserves the Hodge decomposition after every scalar conjugation.
Write \(\overline W\) for scalar complex conjugation.  A \emph{pair
tensor} between \(W,W'\) is a tensor in \(W\otimes W'\) inducing an
isomorphism \(W^*\to W'\), whose every scalar conjugate has type
\((1,1)\).  Symmetry and alternation refer to ordinary tensor
interchange, before the cohomological Koszul sign.  For \(r=\dim W\), put
\[
 D(W)=\bigwedge^rW\subset W^{\otimes r}\subset H^r(X^r,k),
 \qquad D(0)=k
\]
on separate tensor slots, with \(D(0)\) carried by a point.

For a polarizable rational weight-one Hodge structure \(H\) of types
\((1,0),(0,1)\), write \(\operatorname{Hdg}(H)\) for the smallest
rational algebraic subgroup of \(\operatorname{GL}(H)\) containing the
Hodge circle.  It is connected reductive and lies in the symplectic
group of a polarization.  A tensor in \(H^{\otimes 2a}\) is
\emph{balanced} when it has type \((a,a)\).

\begin{companionresults}[Algebraic operations and degree-one tensors]
\label{comp:tools}
\leavevmode
\begin{enumerate}[label=\textup{(\roman*)}]
\item \emph{CM operations.}  CM supply is preserved by external
products, permutations, cup products, rational algebraic
correspondences, \(k\)-linear combinations of such correspondences,
and scalar conjugation.  Every rational Hodge class in a CM-supplied
subspace is algebraic.
\cite[Lemma~2.1]{OpenAIAbelianPowers}.

\item \emph{Admissible operations.}  Every admissible map is induced on
the indicated summands by a \(k\)-linear combination of rational
algebraic correspondences.  Pair tensors are algebraic over \(k\), and
contraction by the inverse pairing is an algebraic cohomological
operation.  If \(W\) is admissible, then \(\overline W\) is admissible
and paired with \(W\).  Alternation and these operations factor
determinant lines across admissible direct sums and contract paired
determinants nontrivially; an alternating self-pair supplies its top
exterior line.
\cite[Lemma~2.2 and its determinant consequences]{OpenAIAbelianPowers}.

\item \emph{Tensor stabilizer.}  Inside the symplectic group of \(H\),
\(\operatorname{Hdg}(H)\) is the pointwise stabilizer of all rational
balanced tensors in \(H^{\otimes 2a}\), for all \(a\geq0\).
\cite[Lemma~2.3]{OpenAIAbelianPowers}.

\item \emph{Toric summands.}  If \(H\subset H^1(A,\Q)\) is a rational
Hodge summand of a complex abelian variety and \(\operatorname{Hdg}(H)\)
is a torus, then \(H\simeq H^1(B,\Q)\) for a CM isogeny factor \(B\) of
\(A\).  Its scalar extension decomposes into CM embedding lines, with
repeated occurrences distinguished.
\cite[Lemma~2.4]{OpenAIAbelianPowers}.
\end{enumerate}
\end{companionresults}

\begin{companionresults}[Saturated determinant criterion]
\label{comp:saturated}
Let \(A\) be a complex
abelian variety, \(H=H^1(A,\Q)\), \(J=\operatorname{Hdg}(H)\), and
\(R=J^{\mathrm{der}}\).  The following data are called saturated:
\begin{enumerate}[label=\textnormal{(\alph*)}]
\item a rational Hodge summand \(H_0\subset H\) with toric Hodge group,
together with a decomposition of \(H_{0,k}\) into CM embedding lines,
with repeated occurrences distinguished;
\item a decomposition of the complement of \(H_{0,k}\) into blocks of
nonzero admissible spaces, each admissibly identified with an actual
representative \(V_i\) carrying an alternating or symmetric self-pair,
or, in a linear block with \(\dim V_i\geq2\), with \(V_i\) or
\(\overline V_i\), paired with each other; in every block \(V_i\) is
chosen as an actual summand of that block;
\item containment in the image of \(R_k\) of the product
\[
 \prod_{i\in I_{\mathrm{Sp}}}\operatorname{Sp}(V_i)
 \times\prod_{i\in I_{\mathrm{SO}}}\operatorname{SO}(V_i)
 \times\prod_{i\in I_{\mathrm{SL}}}\operatorname{SL}(V_i),
\]
acting standardly on its own copies, dually on the opposite linear
copies, and trivially on all other blocks and on \(H_{0,k}\);
\item CM supply for every singleton \(D(V_i)\) in an orthogonal or
linear block, without a balance condition.
\end{enumerate}
The symplectic and orthogonal forms in \textnormal{(c)} are inverse to
the specified pair tensors.  The pair in a linear block identifies
\(\overline V_i\) with \(V_i^*(-1)\) as a Hodge-graded space and with
the dual as a group representation.  Saturated data imply that every
rational balanced tensor is algebraic on its separate \(A\)-slots and
that the rational Hodge conjecture holds on every self-power of \(A\).
For every CM abelian variety \(M\), the induced data on \(A\times M\)
remain saturated and give the same conclusions on all its powers.
This is the saturated clause of
\cite[Theorem~2.5]{OpenAIAbelianPowers}.
\end{companionresults}

Its separate determinant supply is what makes the cover locus
independent of the CM companion.  The exact all-period
determinant-balance criterion and its proof remain in the companion.

\begin{companionresults}[Correspondence transfer]
\label{comp:transfer}
Let \(\mathscr S\) be a
collection of smooth projective varieties such that the rational Hodge
conjecture holds on every finite product of its members.  Suppose that,
for each smooth projective \(Y_i\) in a finite list and each degree
\(q\), the group \(H^q(Y_i,\Q)\) is spanned by maps
\[
 \Gamma_*:H^{q-2j}(S,\Q)\longrightarrow H^q(Y_i,\Q),
 \qquad
 \Gamma\in\CH^{\dim S+j}(S\times Y_i)_\Q,\quad S\in\mathscr S .
\]
Then the rational Hodge conjecture holds on every finite product of
the \(Y_i\), with repetitions and with additional members of
\(\mathscr S\).  The same holds for restrictions of maps from
\(k\)-linear combinations of rational algebraic correspondences if
their images span the full target degree after scalar extension.
\cite[Lemma~2.6]{OpenAIAbelianPowers}.
\end{companionresults}
\section{CM sources for character determinants}
\label{p4:sec:determinants}

For an abelian cover, the invariant-tensor theorem will require the
top alternating line of each character space in first cohomology.
The construction below realizes that line on a finite quotient of a
power of the curve.  The relevant character cohomology is detected
on a distinguished Abel fiber; in the branched case, bundles over
Fermat hypersurfaces dominate that fiber, while in the unramified case
it is a union of projective spaces.  We first give the precise
algebraic supply of those hypersurfaces.  The same supply will
later feed the Gale correspondence in Section~\ref{sec:gale}.

\subsection{Fermat cohomology from CM abelian varieties}

\begin{lemma}[Finite quotient cohomology]\label{lem:quotient}
If a finite group \(G\) acts algebraically on a projective variety
\(V\), the quotient map \(q:V\to V/G\) induces
\begin{equation}
 \label{p4:finite-quotient}
 q^*:H^i(V/G,\Q)\xrightarrow{\;\sim\;}H^i(V,\Q)^G.
\end{equation}
The assertion remains true after extending coefficients.  Pullback is
a morphism of mixed Hodge structures; if \(V\) is smooth and projective,
the cohomology of \(V/G\) is pure of the indicated degree.
\end{lemma}
\begin{proof}
The fibers of \(q\) are finite, so proper base change gives
\(R^j q_*\Q=0\) for \(j>0\).  Each fiber is one \(G\)-orbit, hence
\((q_*\Q)^G=\Q\).  Averaging makes invariants exact, and Leray gives
\eqref{p4:finite-quotient}, also after extending coefficients.
Pullback respects mixed Hodge structures.  If \(V\) is smooth and
projective, strictness of the weight filtration and the displayed
injection into \(H^i(V,\Q)\) give purity.
\end{proof}

\begin{lemma}[Bundle and blowup supply]
\label{p4:bundle-blowup-supply}
If the full cohomology of a smooth projective variety $S$ is
CM-supplied, the same holds for every projective bundle over $S$.
If $Z\subset X$ is a nonempty smooth closed subvariety of pure
codimension $c\ge1$ in a smooth projective variety, and the full
cohomologies of $X$ and $Z$ are CM-supplied, the same holds for
$\operatorname{Bl}_Z X$.
In both assertions finite correspondence supplies for the inputs
give finite supplies for the output in each degree.
\end{lemma}

\begin{proof}
For a projective bundle $p:\mathbb P(E)\to S$ of relative
dimension $e-1$, put $\xi=c_1(\mathcal O_{\mathbb P(E)}(1))$.
The projective-bundle formula is
\[
 H^q(\mathbb P(E),\Q)
 =\bigoplus_{i=0}^{e-1}
 \xi^i p^*H^{q-2i}(S,\Q).
\]
The powers of $\xi$ restrict to a basis on each fiber, which proves
the formula by Leray--Hirsch.  Every displayed map is algebraic.

For a blowup $b:\widetilde X\to X$ of a center of codimension
$c$, let $j:E_Z\hookrightarrow\widetilde X$ be the exceptional
divisor and $p:E_Z=\mathbb P(N_{Z/X})\to Z$.  With
$\xi=c_1(\mathcal O_{E_Z}(1))$, the blowup formula is
\[
 H^q(\widetilde X,\Q)=b^*H^q(X,\Q)
 \oplus\bigoplus_{i=1}^{c-1}
 j_*\bigl(\xi^{i-1}p^*H^{q-2i}(Z,\Q)\bigr).
\]
For completeness, the map of relative pairs
$(\widetilde X,E_Z)\to(X,Z)$ induces an isomorphism in relative
cohomology, since it is an isomorphism off the closed subsets and
the quotient spaces obtained by collapsing those subsets agree.
On $E_Z$, the projective-bundle formula and
$j^*j_*\eta=-\xi\eta$ show that subtracting the displayed
exceptional terms leaves a class whose restriction comes from
$Z$.  The long exact sequences of the pairs then show that the
remaining class comes from $b^*H^q(X)$.  The sum is direct:
$b_*b^*=\operatorname{id}$, $b_*$ kills these exceptional
terms, and their restrictions have independent positive powers
of $\xi$ on $E_Z$.  All maps in the formula are algebraic.
Companion results~\ref{comp:tools}\textup{(i)} prove both supply assertions.
The formulas contain only finitely many summands, proving the
last assertion.
\end{proof}

\begin{proposition}[Fermat supply]
\label{p4:fermat-supply}
For integers $n\ge2$ and $a\ge0$, let
\[
 F_n^a=\{x_0^n+\cdots+x_{a+1}^n=0\}
       \subset\mathbb P^{a+1}_\C.
\]
For every degree $q$, there are finitely many CM abelian varieties
$B_\nu$, integers $j_\nu$, and rational algebraic correspondences
$\Gamma_\nu\in\CH^{\dim B_\nu+j_\nu}(B_\nu\times F_n^a)_\Q$
such that the sum of their maps
\[
 \bigoplus_\nu H^{q-2j_\nu}(B_\nu,\Q)(-j_\nu)
 \xrightarrow{\;\sum_\nu(\Gamma_\nu)_*\;}
 H^q(F_n^a,\Q)
\]
is surjective.  One may take every $B_\nu$ to be a power of the
CM abelian variety $\operatorname{Jac}(F_n^1)$, allowing its
zeroth power, a point.  In dimension zero the finite reduced
scheme $F_n^0$ is supplied component by component by points.
\end{proposition}

\begin{proof}
The Fermat hypersurfaces are smooth.  The dimension-zero assertion
is immediate.  The plane curve $F_n^1$ is connected: distinct
positive-degree components in $\mathbb P^2$ would meet by
B\'ezout's theorem, making their union singular.  For this curve,
the coordinate power map
is the quotient by
$G=(\mu_n)^3/\mu_n$ onto the line
$z_0+z_1+z_2=0$ in $\mathbb P^2$, branched only at its three
coordinate points.  The invariant part of first cohomology vanishes
by \eqref{p4:finite-quotient}.  For a nontrivial character
$\chi$ of $G$, its summand of the finite pushforward of the
constant sheaf is a rank-one local system off the branch points
where $\chi$ is nontrivial on inertia, and has zero stalk at
those points.  At a point where $\chi$ kills inertia it extends
across the point, as one sees by passing to the quotient by
$\ker\chi$.  If there are $r\le3$ active points, its Euler
characteristic is $2-r$.  The connected curve has no nontrivial
character in degrees zero or two, so
\[
 \dim_k H^1(F_n^1,k)_\chi=r-2\le1.
\]
Here a character occurring in first cohomology necessarily has
$r\ge2$.  The Hodge group of the Jacobian commutes with the
algebraic deck action.  Over $k$ its centralizer on first
cohomology is a product of multiplicative groups, since every
occurring character has multiplicity one.  The connected Hodge
group is therefore a torus.  Thus $\operatorname{Jac}(F_n^1)$
is CM.  The Abel map supplies the curve's first cohomology from
this Jacobian, and point correspondences supply degrees zero and
two.

We recall the explicit domination associated with Shioda and
Katsura \cite{SK}, including its resolution so that
the maps used for supply are algebraic.  For $a,b\ge1$, choose
$\varepsilon\in\C$ with $\varepsilon^n=-1$.  There is a
rational map
\begin{equation}
 \label{p4:eq:fermat-map}
 \begin{split}
 F_n^a\times F_n^b&\dashrightarrow F_n^{a+b},\\
 (x,y)&\longmapsto
 \bigl((x_i y_{b+1})_{0\le i\le a},
       (\varepsilon x_{a+1}y_j)_{0\le j\le b}\bigr).
 \end{split}
\end{equation}
Summing the $n$th powers of the two blocks gives the target
equation.  To see dominance, take a target point whose first
block has nonzero power sum.  Rescale its projective coordinates
so that this sum is $-1$; the second block then has power sum
$1$.  Set $x_{a+1}=y_{b+1}=1$, use the first block for
$(x_0,\ldots,x_a)$ and $\varepsilon^{-1}$ times the second
for $(y_0,\ldots,y_b)$.  These are preimages on the two Fermat
hypersurfaces.  Source and target have equal dimension, so the
map is generically finite.

The base locus is the smooth codimension-two subvariety
\[
 Z=\{x_{a+1}=y_{b+1}=0\}
   \simeq F_n^{a-1}\times F_n^{b-1}.
\]
Near each of its points some $x_i$ and some $y_j$ in the
displayed blocks are units in projective coordinates.  The base
ideal of \eqref{p4:eq:fermat-map} is consequently exactly the
ideal generated by $x_{a+1},y_{b+1}$.  Blowing up $Z$
removes this ideal and extends the map to a morphism
$f:\widetilde X\to F_n^{a+b}$ from a smooth projective variety.

Induct on the target dimension, using $a=1$, $b=d-1$ in
dimension $d\ge2$.  The product $F_n^1\times F_n^{d-1}$
and the center $F_n^0\times F_n^{d-2}$ are supplied by
finitely many correspondence maps from powers of $\operatorname{Jac}(F_n^1)$,
by the induction hypothesis and external products.
The center may be disconnected; the blowup formula and supply
apply to its finitely many smooth components.
Lemma~\ref{p4:bundle-blowup-supply} gives the same supply for
$\widetilde X$.  If $e>0$ is the generic degree of $f$,
the projection formula gives $f_*f^*=e\,\operatorname{id}$
on rational cohomology.  Thus $f_*$ is surjective, and composing
the supply maps with it supplies the full target cohomology.
Every source remains a power of the same Jacobian, and each
stage uses finitely many correspondence maps in each degree.
\end{proof}

\subsection{Character determinants on curves}

Every nontrivial character of an abelian deck group becomes faithful
on the cyclic quotient by its kernel.  We first construct the supply
for a cyclic cover, then
transfer it to the original curve in
Corollary~\ref{p4:abelian-determinants}.

We will use the following elementary vanishing.  For a rank-one
$k$-local system on a circle with monodromy $a\ne1$, the cellular
cochain complex is $k\xrightarrow{a-1}k$, hence is acyclic.
The cellular complex for a product of circles is the tensor
product of the circle complexes.  It follows that a rank-one
local system with at least one nontrivial monodromy has zero
cohomology on $(\C^\times)^s$, and also on its product with
an affine space.  Its compactly supported cohomology vanishes
as well, by Poincar\'e duality with the dual local system.
The same cellular argument gives ordinary cohomology vanishing
for a nontrivial rank-one local system on a compact real torus.

\begin{lemma}[Faithful cyclic determinants]
\label{p4:cyclic-determinant}
Let $q:C\to B$ be a connected cyclic cover of smooth connected
projective complex curves, with group $\Delta$ of order $m\ge2$.
Let $\chi:\Delta\to k^\times$ be faithful and put
$V_\chi=H^1(C,k)_\chi$, using inverse pullbacks for the group
action.  Then $D(V_\chi)$ is CM-supplied on its product of
curve slots.  The assertion includes $V_\chi=0$.

More precisely, the CM sources in the defining correspondence
span may all be taken to be powers of the fixed CM abelian variety
$\operatorname{Jac}(F_m^1)$, allowing its zeroth power, a point.
\end{lemma}

\begin{proof}
The group-algebra character idempotents and all their scalar
conjugates preserve the Hodge decomposition, so $V_\chi$ is
admissible and its determinant has the meaning fixed in
Section~\ref{sec:shared}.  We use only the closure assertions
of Companion results~\ref{comp:tools}\textup{(i)}, not the Hodge-class conclusion.

Let $g$ be the genus of $B$ and let
$D=t_1+\cdots+t_r$ be its reduced branch divisor.
The images of the branch meridians have product one in $\Delta$,
by the relation in the punctured surface group and commutativity
of $\Delta$.  Thus $r=0$ or $r\ge2$.  The $\chi$-part of
$q_*k_C$ is a rank-one local system on $B\setminus D$, with
zero stalk at every point of $D$, because $\chi$ is faithful
on each nontrivial inertia group.  Its Euler characteristic is
$2-2g-r$.  Since $C$ is connected, the nontrivial character
does not occur in $H^0(C,k)$ or $H^2(C,k)$.  Hence
\begin{equation}
 \label{p4:eq:cyclic-rank}
 n:=\dim_k V_\chi=2g+r-2.
\end{equation}
There is no nontrivial connected unramified cover of a genus-zero
base, whose fundamental group is trivial.  The numerical cases
$n=0$ are therefore $g=0,r=2$ and $g=1,r=0$; in either case
the determinant is the unit supplied by a point.  We assume
$n>0$ from now on.

\paragraph{The determinant on a finite quotient.}
Set $X=C^n$ and $G=\Delta^n\rtimes\mathfrak S_n$.
Multiplication of the cyclic coordinates is a homomorphism
$G\to\Delta$; let $N$ be its kernel.  The finite quotients are
\begin{equation}
 \label{p4:eq:cyclic-quotients}
 Q:X\longrightarrow Z=X/N,\qquad
 \pi:Z\longrightarrow Z/\Delta
       =B^{(n)}:=\operatorname{Sym}^n B .
\end{equation}
The symmetric power is smooth: in a small disk about each
distinct point of a divisor, elementary symmetric functions
are coordinates for its cluster of points.

By \eqref{p4:finite-quotient}, the pullback of
$H^n(Z,k)_\chi$ is the subspace of $H^n(X,k)$ on which
$(\delta_1,\ldots,\delta_n;\tau)$ acts by
$\chi(\delta_1\cdots\delta_n)$.  On each curve factor this
nontrivial character excludes degrees zero and two.  The only
possible K\"unneth summand is therefore $V_\chi^{\otimes n}$.
A geometric permutation of these degree-one factors includes
its Koszul sign, so invariance under $\mathfrak S_n$ means
ordinary antisymmetry.  Consequently
\begin{equation}
 \label{p4:eq:cyclic-determinant-line}
 Q^*H^n(Z,k)_\chi=D(V_\chi).
\end{equation}
In particular $H^n(Z,k)_\chi$ is a line.

Put $U=(B\setminus D)^{(n)}$ and
$\mathcal L=(\pi_*k_Z)_\chi$.  At a divisor meeting a branch
point, the stabilizer of a lifted tuple contains the inertia
acting in that coordinate.  Its image under the product
character is nontrivial.  The $\chi$-part of the permutation
representation on that fiber is therefore zero.

Over $U$, the map $\pi$ is a $\Delta$-torsor, including along
the diagonals.  Indeed, suppose that a deck tuple followed by a
permutation fixes a lifted tuple.  Around each cycle of the
permutation, the product of the deck displacements fixes one
point of the unramified cover and hence is the identity.  Their
total product is therefore the identity, so the residual
$\Delta$-action is free.  To check local triviality at repeated
points, choose disjoint disks about the distinct points of the
base divisor and label the unramified sheets over each disk.
Quotienting the sheet labels by $N$ retains only their total
product, while the coordinate permutations give the symmetric
powers of the disks.  This yields $m$ local copies of the
base neighborhood.  Thus $\mathcal L$ is extension by zero
of a rank-one local system on $U$.

Its monodromy can be read from moving divisors.  Follow the
$n$ points along a loop and multiply the deck displacements
of their lifted paths; their order at the endpoint may be
permuted.  Closing these paths to a fixed base point cancels
the auxiliary paths because $\Delta$ is abelian.  The product
is the original cover monodromy on the sum of the moving
paths in $H_1(B\setminus D,\mathbf Z)$, with multiplicities.
One may first make this calculation off the diagonals and
perturb a loop to avoid them; the local triviality just proved
gives the same rule on all of $U$.

For both cases below, let
\[
 u:B^{(n)}\longrightarrow\operatorname{Pic}^n B
\]
be the Abel map, sending an effective divisor to its line bundle.
The fiber over $L$ is its complete linear system
$|L|=\mathbb P H^0(B,L)$, where projective space parametrizes
lines.  In both cases $\deg K_B(D)=n$.  Put
\[
 L_0=K_B(D),\qquad
 T=(\pi^{-1}(|L_0|))_{\mathrm{red}}.
\]
Schoen's construction also uses this distinguished Abel fiber
\cite[Section~2, pp.~12--13]{Schoen}.
We will show that restriction to $T$ detects $H^n(Z,k)_\chi$
and construct a proper surjection onto $T$ from smooth projective
varieties with CM-supplied cohomology.  Detection maps the determinant
into source cohomology.  After the two geometric cases, purity and
Poincar\'e duality will turn this detection into supply in the required
direction, through a correspondence obtained from the fiber product
with $C^n\to Z$.

\paragraph{Localization in the branched case.}
Suppose $r>0$.  Then $n\ge2g$.  For $L\ne L_0$, the evaluations
\begin{equation}
 \label{p4:eq:cyclic-evaluations}
 H^0(B,L)\longrightarrow\bigoplus_{i=1}^r L|_{t_i}
\end{equation}
are surjective.  In fact, $H^1(B,L)=0$ since
$\deg L=n>2g-2$.  By the evaluation sequence and Serre duality,
the obstruction is dual to
$H^0(B,K_B\otimes L^{-1}(D))$.  This is the space of sections
of a nontrivial degree-zero line bundle and is zero.
Riemann--Roch gives kernel dimension $g-1$.
If an $L\ne L_0$ exists then $g\ge1$, since
$\operatorname{Pic}^n\mathbb P^1$ is a point.
Coordinates adapted to the independent evaluations give
\begin{equation}
 \label{p4:eq:cyclic-generic-complement}
 |L|\cap U\simeq\mathbb A^{g-1}\times(\C^\times)^{r-1}.
\end{equation}

The restriction of $\mathcal L$ to this complement is nontrivial.
Choose a section on one evaluation hyperplane and no other,
and take a transverse analytic disk through it in $|L|$.
The chosen evaluation has a simple zero in the disk parameter.
Even if the section has a multiple zero at that branch point,
Weierstrass preparation says that the product of the nearby
root coordinates is that evaluation times a nonvanishing
function.  On the boundary of the disk the sum of the winding
numbers of those roots about the branch point is therefore one.
All other roots remain in disks away from the branch set.
The moving-divisor rule identifies the monodromy with the
nontrivial inertia character there.  Collisions away from
the branch set can be perturbed apart, by local triviality
across the diagonals.

The torus vanishing above now makes the compactly supported
cohomology of \eqref{p4:eq:cyclic-generic-complement} with
these coefficients zero.  Because $\mathcal L$ is extension
by zero and $|L|$ is compact, this is
$H^*(|L|,\mathcal L|_{|L|})$.  Proper base change shows that
each $R^j u_*\mathcal L$ is supported at $L_0$.  A sheaf
supported at one closed point has no higher cohomology.
The Leray edge map, followed by proper base change for $\pi$,
therefore gives the restriction isomorphism
\begin{equation}
 \label{p4:eq:cyclic-exceptional-restriction}
 H^n(Z,k)_\chi\xrightarrow{\;\sim\;}H^n(T,k)_\chi.
\end{equation}
When $g=0$ there is no other Picard fiber, so the same
conclusion is immediate.  Reduction does not change the
underlying cohomology.

\paragraph{Fermat sources over the exceptional fiber.}
We next construct a proper surjection onto $T$ from smooth
projective varieties whose full cohomology is CM-supplied.
The residue sequence and the residue theorem give an exact
sequence
\[
 0\longrightarrow G_0:=H^0(B,K_B)
 \longrightarrow H^0(B,K_B(D))
 \longrightarrow
 R:=\{(w_1,\ldots,w_r)\in\C^r:\textstyle\sum_iw_i=0\}
 \longrightarrow0 .
\]
The surjectivity on the right also follows by Riemann--Roch.
Choose a vector-space splitting, so
$|L_0|=\mathbb P(G_0\oplus R)$.  Its evaluation hyperplane
at $t_i$ is $w_i=0$: a meromorphic differential with at most
a simple pole there vanishes at $t_i$ as a section of $K_B(D)$
exactly when its residue is zero.  These hyperplanes are
distinct for $r\ge3$; the two hyperplanes coincide for $r=2$.

The finite reduced map $T\to|L_0|$ is still a quotient by
$\Delta$.  Averaging shows that invariants commute with
quotienting the ideal defining $|L_0|$, and taking reductions
leaves its reduced quotient unchanged.  Every component of
$T$ dominates $|L_0|$.  At least one does, since the finitely
many closed images of the components cover an irreducible base.
Its image is then the whole base, so it meets every fiber.
Each fiber is a single $\Delta$-orbit, and its translates
therefore cover $T$.  Every component is one of those translates.

Let $F=\C(|L_0|)$.  Over $F$ the cover is a possibly disconnected
$\mu_m$-torsor, after identifying $\Delta$ with $\mu_m$.
Kummer theory writes its algebra as $F[z]/(z^m-f)$ for some
$f\in F^\times$.  It is unramified away from the distinct
evaluation hyperplanes.  Thus the valuation of $f$ at every
other prime divisor is divisible by $m$: a Kummer class extending
over that divisor's discrete valuation ring is represented there
by a unit times an $m$th power.

Write the distinct hyperplanes as $H_i=\{\ell_i=0\}$,
$1\le i\le s$, with $\ell_i$ chosen among the residue coordinates.
For $i\ge2$, choose an integer $a_i$ representing the valuation
of $f$ at $H_i$ modulo $m$.  The degree of a principal divisor
is zero.  Since each $H_i$ has degree one and every other
valuation is divisible by $m$, the valuation at $H_1$ is
$-\sum_{i\ge2}a_i$ modulo $m$.  Hence
\[
 \operatorname{div}\left(
 f\Big/\prod_{i=2}^s(\ell_i/\ell_1)^{a_i}\right)=mE
\]
for an integral divisor $E$ on $|L_0|$.  The Picard group of
projective space is torsion-free, so $E$ is principal.
It follows that
\begin{equation}
 \label{p4:eq:cyclic-kummer-class}
 f=c h^m\prod_{i=2}^s(\ell_i/\ell_1)^{a_i},
 \qquad c\in\C^\times,\quad h\in F^\times.
\end{equation}
This includes all prime divisors, not only those in an affine
chart.  When $r=2$ we have $s=1$, so the product is empty and
the generic torsor splits.

The coordinate power map
\[
 s_m:F_m^{r-2}\longrightarrow\mathbb P R,\qquad
 [u_1:\cdots:u_r]\longmapsto[u_1^m:\cdots:u_r^m]
\]
is finite and surjective.  Form the smooth projective bundle
\begin{equation}
 \label{p4:eq:cyclic-fermat-bundle}
 Y=\mathbb P_{F_m^{r-2}}
 \left((G_0\otimes\mathcal O)\oplus
              s_m^*\mathcal O_{\mathbb P R}(-1)\right).
\end{equation}
The tautological subline of $R$ defines a morphism
$Y\to\mathbb P(G_0\oplus R)=|L_0|$.  On each component this
map is dominant and generically finite: away from zero
$R$-coordinate its fibers are the finite fibers of $s_m$.
The pullbacks of the hyperplane ratios in
\eqref{p4:eq:cyclic-kummer-class} are $(u_i/u_1)^m$.
The scalar $c$ also has an $m$th root in $\C$.  The generic
torsor algebra therefore splits over the function field of
each component of $Y$.  Evaluating at its roots gives rational
lifts to every prescribed component of $T$.

Each rational lift extends to a morphism.  Indeed, for a normal
integral variety mapping to a base, a rational lift to a scheme
finite over that base sends each element of the finite affine
algebra to a rational function integral over every local ring
of the source.  Normality makes it regular; the resulting local
lifts glue.  The image here is closed by properness and contains
the generic point of the prescribed component of $T$, so the
lift surjects onto it.  Taking disjoint copies of the required
components of $Y$ gives a proper surjection
\begin{equation}
 \label{p4:eq:cyclic-smooth-surjection}
 Y'\longrightarrow T
\end{equation}
from smooth projective components of dimension $g+r-2$.
Proposition~\ref{p4:fermat-supply} and
Lemma~\ref{p4:bundle-blowup-supply} supply their full cohomology.
The vector space $G_0$ in \eqref{p4:eq:cyclic-fermat-bundle}
is constant, so this construction uses no cohomology of $B$.
For $r=2$, $F_m^0$ is a finite set and $Y'$ is a union of
projective spaces, as needed.

\paragraph{Localization and supply in the unramified case.}
Suppose instead that $r=0$.  Since $n>0$,
\eqref{p4:eq:cyclic-rank} gives $g\ge2$ and $n=2g-2$.
The character of the cover factors through $H_1(B,\mathbf Z)$,
the period lattice of $\operatorname{Jac}(B)$.  It defines a
nontrivial rank-one local system $\mathcal E$ on
$\operatorname{Pic}^n B$, after choosing an origin on that
torsor.  The moving-divisor rule gives $\mathcal L=u^*\mathcal E$.

Riemann--Roch and Serre duality give
\[
 h^0(B,L)=g-1\quad(L\ne K_B),\qquad h^0(B,K_B)=g.
\]
Thus the fibers of $u$ are $\mathbb P^{g-2}$ except for the
fiber $\mathbb P^{g-1}$ over $K_B$.  Choose an ample line
bundle on $B^{(n)}$.  Its $j$th Chern power restricts to a
nonzero generator in degree $2j$ on every fiber for
$0\le j\le g-2$.  It defines, by proper base change, a
stalkwise isomorphism $k\simeq R^{2j}u_*k$ in each such
degree.  All odd direct images vanish, and the only remaining
one is $R^{2g-2}u_*k$, a rank-one skyscraper at $K_B$.

By the projection formula the direct images for $u^*\mathcal E$
are these sheaves tensored with $\mathcal E$.  The cohomology
of $\mathcal E$ on the compact real torus underlying
$\operatorname{Pic}^n B$ vanishes by the calculation at the
start of the subsection.  Leray leaves only the skyscraper
term.  Here $L_0=K_B$, so the Leray edge map and proper base change
give the same restriction isomorphism
\eqref{p4:eq:cyclic-exceptional-restriction}.
The torsor on
$|K_B|=\mathbb P^{g-1}$ splits, since this projective space
is simply connected.  Hence $T$ itself is a disjoint union
of projective spaces of dimension $g-1$.  We may take $Y'=T$;
its full cohomology is supplied by points.

\paragraph{Detection on a smooth source and return to the curve product.}
In either case we have a proper surjection
$Y'\to T$ from smooth projective components whose full
cohomology is CM-supplied.  We use the following precise
consequence of proper cohomological descent:
if $T$ is proper and $Y'\to T$ is a proper surjection with
$Y'$ smooth, then
\begin{equation}
 \label{p4:eq:cyclic-weight-kernel}
 \ker\big(H^n(T,\Q)\longrightarrow H^n(Y',\Q)\big)
       =W_{n-1}H^n(T,\Q).
\end{equation}
This is Deligne's Proposition~8.2.5
\cite[Proposition~8.2.5, p.~39]{Deligne1974}: its proof
injects the weight-$n$ quotient into $H^n(Y',\Q)$.
Properness of $T$ bounds its weights above by $n$, and the
proper smooth $Y'$ has pure degree-$n$ cohomology, giving the
displayed kernel equality.  The proposition requires neither
$T$ irreducible nor $Y'$ connected, so applies here.

By \eqref{p4:finite-quotient}, $H^n(Z,\Q)$ is pure of weight
$n$, since it injects into $H^n(X,\Q)$ and $X=C^n$ is smooth
and projective.  Restriction to $T$ is a morphism of mixed
Hodge structures.  Strictness therefore gives
\[
 \operatorname{im}\big(H^n(Z,\Q)\to H^n(T,\Q)\big)
       \cap W_{n-1}H^n(T,\Q)=0.
\]
Extend scalars to $k$ and take a nonzero
$\alpha\in H^n(Z,k)_\chi$.  Its restriction is nonzero by
\eqref{p4:eq:cyclic-exceptional-restriction}.  It survives
on some component $Y_0$ of $Y'$ by
\eqref{p4:eq:cyclic-weight-kernel}.  Write
$p:Y_0\to T\to Z$ and $d=\dim Y_0$.  Poincar\'e duality
provides
\[
 \beta\in H^{2d-n}(Y_0,k),\qquad
 \int_{Y_0}p^*\alpha\smile\beta\ne0.
\]
Here $(d,2d-n)=(g+r-2,r-2)$ in the branched case and
$(d,2d-n)=(g-1,0)$ in the unramified case.  In particular
these are valid degrees.  The class $\beta$ is CM-supplied
because the full cohomology of $Y_0$ is.

Choose an irreducible component of the finite surjection
$Y_0\times_Z X\to Y_0$ which dominates $Y_0$, and resolve it
projectively.  The resulting smooth projective variety
$\widetilde Y$ fits into the commutative square
\[
 \begin{array}{ccc}
  \widetilde Y & \xrightarrow{\ l\ } & X=C^n\\
  {\scriptstyle h}\big\downarrow && \big\downarrow{\scriptstyle Q}\\
  Y_0 & \xrightarrow{\ p\ } & Z,
 \end{array}
\]
where $p$ factors as $Y_0\to T\to Z$, $h$ has positive generic degree
$e$, and $Ql=ph$.  The component $Y_0$ has full CM supply and detects
$p^*\alpha$; the resolving variety $\widetilde Y$ is used only to realize
the return correspondence $l_*h^*$.
The algebraic correspondence
\[
 \Gamma=(h,l)_*[\widetilde Y]\in\CH^n(Y_0\times X)_\Q
\]
induces $l_*h^*$.  Its cohomological degree shift is
$2(n-d)$, so it sends $\beta$ to $H^n(X,k)$.
This image is CM-supplied by Companion results~\ref{comp:tools}\textup{(i)}.
Only the pullback of $\beta$ is used; no supply assertion
about the cohomology of $\widetilde Y$ is required.
Projection formula gives
\begin{equation}
 \label{p4:eq:cyclic-return-pairing}
 \int_X Q^*\alpha\smile\Gamma_*\beta
 =\int_{\widetilde Y}h^*(p^*\alpha\smile\beta)
 =e\int_{Y_0}p^*\alpha\smile\beta\ne0.
\end{equation}

Let $\rho:G\to k^\times$ be
$\rho(\delta_1,\ldots,\delta_n;\tau)
=\chi(\delta_1\cdots\delta_n)$.  The inverse-character
idempotent is
\[
 P_{\rho^{-1}}=\frac1{|G|}\sum_{a\in G}\rho(a)a.
\]
It is a $k$-linear combination of algebraic automorphism
graphs.  The integration pairing is $G$-invariant, while
$Q^*\alpha$ transforms by $\rho$.  Averaging
\eqref{p4:eq:cyclic-return-pairing} therefore gives
\[
 \int_X Q^*\alpha\smile P_{\rho^{-1}}\Gamma_*\beta
 =\int_X Q^*\alpha\smile\Gamma_*\beta\ne0.
\]
The calculation for \eqref{p4:eq:cyclic-determinant-line}
identifies the image of $P_{\rho^{-1}}$ in degree $n$ with
$D(H^1(C,k)_{\chi^{-1}})$; the inverse character has the same
rank $n$ by \eqref{p4:eq:cyclic-rank}.  This determinant line
is thus CM-supplied.  Complex conjugation of coefficient
scalars sends $\chi^{-1}$ to $\chi$, since its values are
roots of unity.  It preserves the rational correspondence
span by Companion results~\ref{comp:tools}\textup{(i)}, proving the assertion
for $D(V_\chi)$.

Finally, in the branched case
Proposition~\ref{p4:fermat-supply} starts the supply in
\eqref{p4:eq:cyclic-fermat-bundle} from powers of
$\operatorname{Jac}(F_m^1)$.  The projective-bundle maps
and all subsequent maps merely compose those rational
correspondences, and coefficient conjugation preserves their
rational image span.  The unramified and zero-rank cases
use a point.  This proves the stated fixed-source refinement.
\end{proof}

\begin{corollary}[Determinants for abelian covers]
\label{p4:abelian-determinants}
Let a finite abelian group $P$ act faithfully on a smooth connected
projective complex curve $C$, and put $B=C/P$.  For every
character $\chi:P\to k^\times$, the line
$D(H^1(C,k)_\chi)$ is CM-supplied.  This includes the zero
character spaces.  For the trivial character the line is
generated by a rational algebraic cycle class.

For $\chi\ne1$, let $r_\chi$ be the number of branch points of
$C/\ker\chi\to B$.  Then
\[
 \dim_k H^1(C,k)_\chi=2g(B)+r_\chi-2.
\]
The number $r_\chi$ counts precisely the original branch points
whose inertia is not killed by $\chi$.

For a fixed $P$, the CM sources for all these lines can be taken
to be powers of one fixed product of Fermat-curve Jacobians
depending only on $P$, not on the complex structure of $C$.
\end{corollary}

\begin{proof}
Suppose first that $\chi\ne1$ and put $K=\ker\chi$.
The quotient $C_\chi=C/K$ is a smooth connected projective
curve; locally the quotient of a smooth curve by its finite
stabilizer has parameter $t^e$.  The residual group $P/K$
is cyclic and $\chi$ is faithful on it.  Its action on
$C_\chi$ is effective.  Indeed, choose a point of $C$ with
trivial $P$-stabilizer, outside the finitely many fixed points
of nonidentity automorphisms.  If a coset acted trivially on
$C/K$, it would send this point into its $K$-orbit, forcing
the coset to be the identity.

Equivariantly for the residual action,
\eqref{p4:finite-quotient} identifies
\[
 H^1(C_\chi,k)_\chi\xrightarrow{\;\sim\;}H^1(C,k)_\chi
\]
by quotient pullback.  A point of $B$ is branched in
$C_\chi\to B$ exactly when its original inertia has nontrivial
image in $P/K$.  Lemma~\ref{p4:cyclic-determinant} applies to
this connected cyclic quotient and gives its determinant supply;
\eqref{p4:eq:cyclic-rank} gives the dimension formula.
The tensor power of the rational graph pullback sends that determinant line to
$D(H^1(C,k)_\chi)$, by Companion results~\ref{comp:tools}\textup{(i)}.
If the character space is zero both lines are the unit.
This includes a quotient that becomes unramified although
the original cover is branched.

For $\chi=1$, quotient pullback identifies the character
space with $H^1(B,k)$.  Choose a point $b\in B$.
The rational algebraic class
\[
 [\Delta_B]-[b\times B]-[B\times b]
\]
is the $H^1(B)\otimes H^1(B)$ component of the diagonal and,
up to the cup-pairing convention, is the inverse symplectic
tensor.  If $g=g(B)$, antisymmetrizing its $g$th external
power gives a nonzero multiple of the top exterior power
of this nondegenerate alternating tensor.  It therefore
generates $D(H^1(B,k))$ by a rational algebraic class.
For $g=0$ this is the unit.  Pullback on the curve slots
proves the assertion.

For the last assertion, let $\mathcal O(P)$ be the finite set
of orders of nontrivial characters of $P$, and set
\[
 M_P=\prod_{m\in\mathcal O(P)}\operatorname{Jac}(F_m^1),
\]
with the empty product interpreted as a point.  The faithful
lemma uses powers of the indicated factor for a character
of order $m$, and the trivial case uses a point.  Each such
factor power is supplied from the corresponding power of
$M_P$ by its algebraic factor inclusion and projection.
The product $M_P$ is CM.  This proves the uniform choice
of sources.
\end{proof}

\section{Abelian covers with full marked variation}
\label{p4:sec:covers}

We now apply the determinant criterion to Jacobians of abelian covers.
The determinant of each nontrivial character space is supplied by its
cyclic quotient, and the invariant determinant by the polarization.
The remaining task is to describe the derived Hodge action:
individual character projections must be standard, and characters on
which the same simple ideal acts must be identified by algebraic
correspondences.  The second point is essential in small base genus.

\subsection{The marked family and its Hodge-generic locus}

Fix a compact connected oriented surface $B_{\mathrm{top}}$ of genus
$h\ge0$, a finite labelled set $D\subset B_{\mathrm{top}}$, a point
$*\notin D$, a finite abelian group $P$, and an epimorphism
\[
 \phi:\pi_1(B_{\mathrm{top}}\setminus D,*)\longrightarrow P.
\]
The meridian of every point in $D$ is required to have nonidentity
image.  The images of the meridians multiply to one, by the surface
relation, and together with the handle images generate $P$.
The epimorphism describes a connected unramified $P$-cover off $D$;
compactifying at a meridian of order $e$ uses the local map $u\mapsto
u^e$.  Denote the resulting compact surface by $C_{\mathrm{top}}$.
The trivial group and the empty branch set are permitted.

Add labelled unbranched auxiliary points if needed so that the marked
base, including $*$, is stable.  On the additionally punctured surface,
extend $\phi$ by sending their meridians to the identity element.
Let $\mathcal T_\phi$ be the connected
Teichm\"uller manifold of the resulting labelled surface.  We vary its
\emph{full} marked complex structure: no branch positions or
base-complex-structure directions are held fixed.  Lifting the charts
off $D$ and filling with the fixed local power maps gives a smooth
proper holomorphic family of compact curves $C_s$, locally on
$\mathcal T_\phi$.  Its first cohomology is a marked polarized
variation on
\[
 H_\Q=H^1(C_{\mathrm{top}},\Q),\qquad
 H_{\mathbb Z}=H^1(C_{\mathrm{top}},\mathbb Z).
\]
As in Section~\ref{sec:shared}, put $k=\Qbar$ and write
$H_K=H_\Q\otimes_\Q K$ for $K=k,\C$.
Auxiliary points do not alter these curves or their cohomology.
They only provide stable marked parameter spaces in the small cases.

Let $\Gamma$ be the stabilizer of the exact epimorphism $\phi$ in the
pure based mapping class group, with all the labels fixed.  It has
finite index because there are finitely many homomorphisms from the
finitely generated surface group to $P$.  After choosing a point above
$*$, every element of $\Gamma$ has a unique lift fixing that point.
The lift commutes with $P$ and acts integrally and symplectically on
$H_\Q$.  It respects the marked variation.  We use inverse pullbacks
for left actions; replacing all actions by their opposites would not
change the algebraic closures.  These are the standard marked-surface
and Dehn-twist conventions; see
\cite[Chapters~2, 3, and~6]{FarbMargalit2012}.

\begin{lemma}[The Hodge-generic locus of the marked family]
\label{p4:generic-tensors}
For $a\ge0$ and $t\in H_\Q^{\otimes 2a}$, let
\[
 Z_t=\{s\in\mathcal T_\phi:
       t_\C\in F^a(H_\C^{\otimes 2a})_s\},\qquad
 \mathcal T_\phi^{\mathrm{Hg}}
 =\mathcal T_\phi\setminus
   \bigcup_{\substack{a\ge0,\ t\in H_\Q^{\otimes 2a}\\
                       Z_t\ne\mathcal T_\phi}}Z_t.
\]
Each $Z_t$ is closed complex analytic.  The complement
$\mathcal T_\phi^{\mathrm{Hg}}$ is nonempty and is the complement of a
countable union of proper closed analytic subsets.  At a point in this
complement, every rational balanced tensor is balanced throughout
$\mathcal T_\phi$.

Here $H_s$ denotes $H_\Q$ with the Hodge structure at $s$.
Let $J_s=\operatorname{Hdg}(H_s)$ and let $L$ be the rational Zariski
closure of $\Gamma$ in the symplectic group of $H_\Q$.  At every
$s\in\mathcal T_\phi^{\mathrm{Hg}}$ one has
$L^\circ\subset J_s$, and the whole $\Gamma$ normalizes $J_s$.
\end{lemma}

\begin{proof}
A rational tensor is real.  In weight $2a$, it has type $(a,a)$
exactly when it belongs to $F^a$: reality then also places it in
$\overline{F^a}$, whose intersection with $F^a$ is the $(a,a)$ part.
Membership in the holomorphic subbundle $F^a$ is the vanishing of the
induced section of its quotient bundle, so $Z_t$ is closed analytic.
There are countably many rational tensors.  A proper analytic subset
of the connected complex manifold is nowhere dense, and Baire's theorem
proves the nonemptiness assertion.  Odd tensor degrees have no balanced
vectors, since the circle element $-1$ acts on them by $-1$.

At a point in the displayed complement, the balanced rational tensors
are exactly those balanced at every marked parameter.  The latter
spaces are preserved by $\Gamma$, because it acts equivariantly on the
whole variation.  By Companion results~\ref{comp:tools}\textup{(iii)}, their pointwise
stabilizer in the symplectic group is $J_s$.  Thus every element of
$\Gamma$ normalizes $J_s$.

An individual globally balanced tensor has finite $\Gamma$-orbit.
We use the finite-action argument of
\cite[Section~4, Lemma~4]{Andre1992FixedPart}.
If its denominator is $q$ in the integral tensor lattice, every
translate still lies in $q^{-1}H_{\mathbb Z}^{\otimes 2a}$, and
its tensor-polarization norm is unchanged by transport.  At a fixed
fiber all its translates lie in the real balanced subspace.  The tensor
polarization, with its fixed weight sign, is definite on that subspace.
A fixed norm sphere there is compact and contains only finitely many
points of this lattice.  The orbit is therefore finite.
The connected algebraic group $L^\circ$ fixes every such tensor, so
Companion results~\ref{comp:tools}\textup{(iii)} give $L^\circ\subset J_s$.
\end{proof}

The definition of $\mathcal T_\phi^{\mathrm{Hg}}$ specifies the entire
exceptional locus used below.  It is a locus in the full marked family;
a proper subfamily inherits the conclusion only after verifying that
these tensor loci pull back properly.

\begin{proposition}[Saturated data for the marked cover family]
\label{p4:cover-saturation}
For every connected topological abelian cover and full marked family
defined above, and every \(s\in\mathcal T_\phi^{\mathrm{Hg}}\), the
Jacobian \(\operatorname{Jac}(C_s)\) admits saturated data in
Companion results~\ref{comp:saturated}.
\end{proposition}

This proposition supplies the tensor input for
Theorem~\ref{p4:abelian-covers}.  Write $C\to B$ for a fiber of the
fixed marked family.  Its character decomposition is
\[
 H_k=\bigoplus_\chi W_\chi,\qquad
 W_\chi=e_\chi H_k,\qquad
 e_\chi=\frac1{|P|}\sum_{g\in P}\chi(g)^{-1}g\in k[P].
\]
Every scalar conjugate of $e_\chi$ is again a deck-character projector,
so it preserves the Hodge decomposition.  Thus the $W_\chi$ are
admissible in the sense of Section~\ref{sec:shared}.  The marking
identifies these spaces and the mapping-class operators across all fibers.

Corollary~\ref{p4:abelian-determinants} already supplies their
determinant lines.  We must establish two properties of their
Hodge-group action.  First, every character space with nontrivial
derived action has standard special linear or symplectic projection.
Second, when one simple Hodge Lie ideal acts on several such spaces,
their standard or dual identifications must preserve the Hodge circle
at every scalar conjugate.  Rank-one twist operators give the
projections; the comparison of spaces will depend on whether the base
has positive genus or is a sphere.

\subsection{A rank-one criterion for standard actions}

\begin{lemma}[Rank-one nilpotents]\label{p4:cover-rank-one}
Let $V$ be a finite-dimensional vector space over $k=\Qbar$, and let a
connected reductive algebraic
subgroup of $\operatorname{GL}(V)$ act irreducibly.  If its Lie algebra contains a
rank-one nilpotent, its derived image is the standard special linear or
standard symplectic group on $V$.
\end{lemma}

\begin{proof}
It is enough to prove the Lie algebra assertion, since the groups in the
conclusion are connected.  Extending scalars from $k$ to $\C$ preserves
irreducibility, semisimplicity, and the rank condition, so we work over
$\C$.  In the symplectic alternative the invariant alternating form
descends to $k$: invariant forms are the kernel of linear equations
over $k$, and the determinant is nonzero on a suitable alternating
form with coefficients in $k$.  Equality of the connected group images then descends under
scalar extension.  The center
acts by scalars.  The Jordan decomposition in an algebraic Lie algebra
shows that a nilpotent element has zero central component and nilpotent
components in its simple ideals.  Write the irreducible representation
of the semisimple image as a tensor product $\bigotimes_i V_i$ of
nontrivial irreducibles for its simple ideals.

Let $N$ be the given rank-one nilpotent.  In particular $N^2=0$.  If two
of its simple-ideal components $N_i,N_j$ were nonzero, choose length-two
terminal Jordan strings in $V_i,V_j$ and kernel vectors in the other
factors.  On the tensor product of the top vectors of these strings,
$N^2$ is $2N_i\otimes N_j\ne0$.  This is impossible.  If just $N_i$ is
nonzero, its rank on $V$ is
$\operatorname{rank}(N_i)\prod_{j\ne i}\dim V_j$.  Rank one therefore forces the
semisimple image to be simple.

Put $N$ in the nilradical of a Borel subalgebra of this simple image.
A generic one-parameter subgroup of its maximal torus, together with
a scalar rescaling, degenerates $N$ to a nonzero root vector $x_\alpha$.
The locus of matrices of rank at most one is closed, so $x_\alpha$ has
rank one.  Every root vector for a root of the same length has rank one,
by Weyl conjugacy.  Restriction to the corresponding root
$\mathfrak{sl}_2$ consequently consists of one standard doublet and
trivial representations: the rank of the raising operator on an
irreducible $\mathfrak{sl}_2$-module of dimension $a$ is $a-1$.
Thus each reflection in a root of that length exchanges exactly two
weights, both of multiplicity one, and fixes every other weight.

Here is the root-system check that finishes the argument.  We use the
usual classification and weight lattices of reduced irreducible root
systems, as in \cite[Sections~23--25 and~30]{Etingof2024Lie}.  The roots
of a specified length form a full-rank simply laced root system, not
necessarily a closed root subsystem of the ambient root system: the whole
system in simply laced type; $D_\ell$ or $A_1^\ell$ in types $B_\ell$
and $C_\ell$; $D_4$ in type $F_4$; and $A_2$ in type $G_2$.  Use the
small-rank conventions $D_3=A_3$ and $D_2=A_1\sqcup A_1$.

The simple reflections of this root system act as distinct
transpositions on the weights.  They are distinct because the weights
span the dual Cartan space.  Orthogonal simple roots give disjoint
transpositions, while adjacent simple roots give transpositions sharing
one endpoint.  A branching vertex of the Dynkin diagram would require
three disjoint transpositions each meeting the same two-element set,
which is impossible.  A connected full-rank path $A_\ell$ therefore
acts on precisely $\ell+1$ moved weights as successive transpositions
on a path.  All other weights, being fixed by this full-rank root system,
are zero.  The moved weights form its standard regular simplex: their
successive differences are the simple roots with a consistent choice
of sign, as follows by pairing with adjacent simple coroots.  In full
type $A_\ell$ the highest weight is thus the first or last fundamental
weight.  Irreducibility gives the standard representation or its dual.
If this path arose instead from $G_2$, $B_3$, or $C_3$, the full Weyl
group contains negation, whereas this simplex, for $\ell\ge2$, is not
stable under negation.  Those cases are excluded.

It remains to consider the orthogonal-axis root system $A_1^\ell$.
Its transpositions are disjoint.  A moved weight for one axis is fixed
by every other axis reflection, so lies on that axis; the doublet
condition makes its two weights $\pm\alpha/2$.  For the long roots
$2e_i$ in $C_\ell$, this gives the weights $\pm e_i$, and the highest
weight is that of the symplectic standard representation.  For the
short roots $e_i$ in $B_\ell$, the vectors $\pm e_i/2$ are not weights
when $\ell\ge2$.  In the rank-two $D_2$ root system the half-roots
$(\pm e_1\pm e_2)/2$ are not weights of $C_2$; in $B_2$ they give its
spin representation, which is the standard representation under
$B_2=C_2$.  Rank one gives the ordinary $\mathfrak{sl}_2$ standard.
These are exactly the alternatives in the statement.
\end{proof}

\subsection{Character cohomology and disk twists}

We next construct the rank-one operators to which the preceding lemma
will apply.  Their formulas also let us compare different characters.

For a character $\chi:P\to k^\times$, the space
$W_\chi$ is $H^1(C,k)_\chi$.  Inverting all character conventions makes no
difference; choose the convention in which $W_\chi$ is computed by the
rank-one character $t=\chi\phi$ below.  The trivial character gives
$W_{\mathbf1}=H^1(B,k)$, by quotient pullback.  For a nontrivial
character, pass to its cyclic quotient and fill every puncture whose
local multiplier is one.  Write the remaining multipliers as
$t_1,\ldots,t_r\ne1$.  Then $r=0$ or $r\ge2$, and
\begin{equation}\label{p4:character-dimension}
 \dim W_\chi=2h+r-2.
\end{equation}

For a nonzero complex character space, \emph{pure} and \emph{mixed}
mean respectively that just one, or both, of its $(1,0)$ and
$(0,1)$ parts are nonzero.

\begin{lemma}[Cocycles and twists]\label{p4:twist-model}
For a nontrivial $t$, the character space is the space of crossed
homomorphisms
\begin{equation}\label{p4:cocycles}
 z(uv)=z(u)+t(u)z(v),
 \qquad z(u)\sim z(u)+(1-t(u))b\quad(b\in k).
\end{equation}
Use generators $a_j,b_j$ for the handles and positively based meridians
$g_i$, with relation $\prod_j[a_j,b_j]\prod_i g_i=1$.  Values of $z$ on
these generators satisfy just the crossed-homomorphism relation on this
word, followed by the one-dimensional coboundary quotient.

Let $\alpha$ be an embedded arc between two active punctures.  A twist
about the boundary of its two-puncture disk acts, in adapted coordinates
and after choosing one of the two twist conventions, by
\begin{equation}\label{p4:pair-twist}
 T_\alpha=1+u_\alpha\ell_\alpha,\qquad
 u_\alpha=(-t_i,1),\qquad
 \ell_\alpha(z)=(1-t_j)z(g_i)+(t_i-1)z(g_j),
\end{equation}
with $u_\alpha$ supported on these two meridians.  Its exceptional
eigenvalue is $t_it_j$.  A different choice of lift can multiply this
formula by a finite-order scalar.

For a nontrivial such transformation, its fixed hyperplane is the
Hermitian orthogonal complement of its image direction.  If a collection
of these directions has connected nonorthogonality graph and their fixed
hyperplanes have zero intersection, the group they generate acts
irreducibly.  These statements apply also to ordinary transvections
arising from simple closed curves with trivial $t$-monodromy.
\end{lemma}

\begin{proof}
The eigenspace on the unbranched cover is the cohomology of the
corresponding rank-one local system on the punctured base, by averaging
over the finite deck group.  Filling an inactive puncture gives ordinary
cohomology on the filled base.  At an active puncture the local complex
is $k\xrightarrow{t_i-1}k$, which is acyclic, so no local term changes
the character cohomology.  The cellular cochain model, or the standard
crossed-homomorphism model of degree-one group cohomology, gives
Equation~\eqref{p4:cocycles}.  Euler characteristic gives
Equation~\eqref{p4:character-dimension}; in the unramified case both degree zero
and degree two vanish for a nontrivial rank-one system.

For the twist, conjugate the two meridians by their product
$\delta=g_ig_j$ and fix the exterior generators.  The identity
\[
 z(fuf^{-1})=(1-t(u))z(f)+t(f)z(u)
\]
gives Equation~\eqref{p4:pair-twist} by direct subtraction.  In
particular $\ell_\alpha(u_\alpha)=t_it_j-1$.  All such disk twists
preserve $\phi$: they act trivially on the abelianization of the
punctured surface group.  For other mapping classes preserving $t$, a
positive power preserves $\phi$, which suffices for all Lie-algebra
statements.

Set $s_i=z(g_i)/(1-t_i)$.  These are the fixed points of the affine
transformations represented by the meridians.  The fixed-hyperplane
condition is $s_i=s_j$ after transport along $\alpha$.  More explicitly,
if $\alpha$ is represented by $p_i^{-1}xp_j$ in terms of based access
paths and a based loop $x$, it is
\begin{equation}\label{p4:fixed-points}
 s_i=z(x)+t(x)s_j.
\end{equation}
This description is unchanged by endpoint windings.  For $r\ge3$ the
image vector in Equation~\eqref{p4:pair-twist} is nonzero modulo
coboundaries, since another active meridian has zero coordinate in it.
When $r=2$, if that vector is a coboundary, all exterior multipliers are
one; the surface relation then also makes $\ell_\alpha$ zero.

The polarization pairs opposite character spaces perfectly and gives
a nondegenerate Hermitian form on $W_{\chi,\C}$, preserved by the
monodromy.  If $T=1+u\ell$ is a nontrivial isometry, every fixed vector
is orthogonal to $(T-1)W=ku$.  Both spaces have codimension one, proving
the fixed-hyperplane assertion.  A common zero intersection of fixed
hyperplanes makes the directions span by nondegeneracy.  A nonzero
invariant subspace contains a vector moved by one transformation, hence
contains its image direction.  Nonorthogonality then propagates this
inclusion along the graph and gives the whole space.

Finally, if a simple closed curve has trivial $t$-monodromy, lift its
annulus to the cyclic cover and twist simultaneously about its lifts.
The ordinary intersection formula restricts on a character space to
rank at most one: the projected lift classes are scalar multiples of
each other.  If evaluation of the cocycle on one closed lift is nonzero,
the intersection-dual projected class is nonzero as well, by the perfect
opposite-character pairing.  The resulting transformation is a
nontrivial rank-one unipotent.  Its logarithm, or a nonzero integral
multiple obtained by taking a power, belongs to the Lie algebra of
connected monodromy.
\end{proof}

\begin{lemma}[Connected irreducibility]\label{p4:connected-irreducibility}
Suppose a subgroup of $\operatorname{GL}(V)$ acts irreducibly and is generated,
up to scalars, by rank-one changes of the identity.  If the Lie algebra
of its Zariski closure contains a rank-one nilpotent, the identity
component of that closure acts irreducibly.
\end{lemma}

\begin{proof}
The unipotent radical of the identity component is normal in the full
closure.  Its common fixed subspace is nonzero and is invariant under
the full group, so it is all of $V$; the radical is trivial in the
given faithful image.  The identity component is therefore reductive.
Its isotypic summands are permuted transitively by the full group.
A rank-one nilpotent can act nontrivially on only one isotypic summand,
and there with multiplicity one; that summand has dimension at least
two.  Transitivity gives the same two properties for every summand.
A scalar multiple of a rank-one change cannot permute two summands of
dimension at least two: its difference from that scalar identity
would have rank at least two.  Thus all the generators preserve every
summand.  Irreducibility forces a single summand, with multiplicity
one, as required.
\end{proof}

\subsection{Positive base genus}

We now assume $h\ge1$.  The handle character can be concentrated on one
generator:
\begin{equation}\label{p4:handle-character}
 t(a_j)=t(b_j)=1\quad(j<h),\qquad t(b_h)=1,\qquad t(a_h)=\eta.
\end{equation}
Indeed lift the character on the handle homology to an integral row,
divide by the greatest common divisor of its entries, and complete its
primitive part to an integral symplectic basis.  The resulting
symplectic basis change is realized by a homeomorphism on the handle
subsurface, disjoint from the puncture disk; see
\cite[Chapter~6]{FarbMargalit2012}.  If all handle multipliers are one,
take $\eta=1$.  This changes the marking, not the topological cover.

\begin{lemma}[Winding arcs]\label{p4:winding-arcs}
Place the active punctures in an ordered disk, away from the handle
subsurface.  For every simple based handle generator $x$, an embedded
arc from puncture $1$ to puncture $2$ can be chosen whose transported
fixed-point condition is Equation~\eqref{p4:fixed-points} with $x$ or
$x^{-1}$.  If $r\ge3$, the arc can avoid the interior of the standard
arc from $2$ to $3$.  Arcs for different generators are chosen
separately; they are not required to be mutually disjoint.

For a winding arc using only the last handle, the simple generators of
the first $h-1$ handles can be represented outside its two-puncture disk.
\end{lemma}

\begin{proof}
Take the based handle bouquet on the handle side of a small vertex disk,
and fan the access paths $p_i$ to the punctures out on its other side.
Choose the first two access paths to avoid the interior of the standard
$2,3$ arc when that arc is present.  Concatenate
$p_1^{-1}xp_2$.  Outside the vertex disk its pieces are embedded and
disjoint.  In the vertex disk there are two germs from the puncture
side and two from the handle side.  Choose the orientation of $x$ so
that the incoming and outgoing germs can be joined in noncrossing
pairs.  Replacing these joins by small smooth arcs removes the vertex
and gives the required embedded arc.  Its regular neighborhood with
the two endpoint disks is a disk; with the $2,3$ arc added it is a
regular neighborhood of an embedded three-vertex tree, also a disk.
This explains why the consecutive-meridian formulas apply to the
winding and standard arcs together.  Auxiliary marks can be avoided
by small perturbations.

Use the standard separated handles for the last assertion.  The winding
arc and its endpoint disk can be placed in the last-handle and puncture
region.  The other handle loops lie outside.  Carrying all these choices
by the homeomorphism used in Equation~\eqref{p4:handle-character}
preserves the assertion.  A change of their based access paths only
conjugates those loops; because their $t$-multipliers are one, zero
cocycle displacement on them remains zero after that change.
\end{proof}

\begin{proposition}[Positive-genus character projections]
\label{p4:positive-projections}
For $h\ge1$, every nonzero nontrivial character space is mixed in Hodge
type at every scalar conjugate.  At every
$s\in\mathcal T_\phi^{\mathrm{Hg}}$ the image
of the derived Hodge group on that space is the full standard special
linear or full standard symplectic group.  The invariant character space,
when nonzero, has Hodge dimensions $(h,h)$ and full standard symplectic
derived action.
\end{proposition}

\begin{proof}
We exhibit an irreducible group of twists and a rank-one nilpotent in
the Lie algebra of its connected closure.  There are three cases.

\emph{No active punctures.}  Here $\eta\ne1$, since $t$ is nontrivial.
The surface relation is $(\eta-1)z(b_h)=0$.  Thus $z(b_h)=0$, and a
coboundary removes $z(a_h)$.  The remaining coordinates are precisely
the ordinary first cohomology coordinates of the first $h-1$ handles.
The twist character is trivial on this subsurface.  Twists along a
symplectic handle basis and along connecting band sums give rank-one
transvections whose directions span with connected nonorthogonality
graph.  Their positive powers preserving $\phi$ have the same
logarithm directions.  They act irreducibly already on the identity
component.  When $h=1$ the character space is zero, as asserted.

\emph{At least three active punctures.}  Start with consecutive arcs
inside the ordered puncture disk.  Their common fixed equations make
all $s_i$ equal.  Adjacent directions pair nontrivially: in a common
three-puncture disk this follows by substituting the vector of one
pair twist into the functional of its neighbor in
Equation~\eqref{p4:pair-twist}; the factors that occur are nonzero
multiples of $t_j$, $1-t_i$, and $1-t_k$.

For each handle generator add the winding arc of
Lemma~\ref{p4:winding-arcs} from $1$ to $2$.  Its direction is linked
to the standard $2,3$ direction by the same calculation in their
three-puncture neighborhood.  In the common fixed space, use a
coboundary to make all $s_i=0$.  The winding equations then force
$z(x)=0$ for every handle generator.  Thus the common fixed space is
zero and the direction graph is connected.  Lemma~\ref{p4:twist-model}
gives irreducibility.

The twist about $b_h$ supplies a rank-one nilpotent.  In compatible
generators it replaces $a_h$ by $a_hb_h$ and fixes the other generators,
so its cocycle displacement is supported on $a_h$ and is
$\eta z(b_h)$.  The coordinate $z(b_h)$ can be chosen freely and the
single surface relation compensated at an active meridian.  The image
vector is not a coboundary, since it vanishes on all active meridians.
This is a nonzero square-zero rank-one operator.  A positive power of
the twist preserves $\phi$; a further power removes any finite deck
scalar.  Adjoin this twist to the irreducible disk-and-winding subgroup.
The enlarged group is still generated, up to scalars, by rank-one
changes of the identity.  Lemma~\ref{p4:connected-irreducibility}
now gives connected irreducibility for its closure.

\emph{Two active punctures.}  Write $t_1=\lambda$,
$t_2=\lambda^{-1}$.  The cocycle relation is
\begin{equation}\label{p4:two-puncture-relation}
 (\eta-1)z(b_h)+(1-\lambda)(s_1-s_2)=0.
\end{equation}
If $\eta=1$, remove the common value $s_1=s_2$ by a coboundary.
All the handle coordinates are then ordinary symplectic coordinates,
and their twists and band sums prove the assertion exactly as in the
first case.

Suppose $\eta\ne1$ and impose the gauge $s_2=0$.  The handle
coordinates are arbitrary and Equation~\eqref{p4:two-puncture-relation}
determines $s_1$.  Take the standard two-puncture arc and the winding
arcs corresponding to $a_h$ and to $a_j,b_j$ for $j<h$.  The standard
fixed equation is $s_1=0$; the others are $s_1=z(x)$, after reversing
$x$ if necessary.  Reversing a trivial-multiplier generator merely
changes its sign, and for $a_h$ it multiplies its coordinate by the
nonzero scalar $-\eta^{-1}$; these changes do not affect the argument.
The common kernel is zero: first $s_1=z(b_h)=0$, then $z(a_h)=0$ and
all remaining handle coordinates vanish.

We verify connectedness of the direction graph.  Before the gauge,
the standard arc direction has zero handle coordinates and
$s_1=s_2=-\lambda/(1-\lambda)\ne0$.  Adding the coboundary with
parameter $\lambda/(1-\lambda)$ puts it in the gauge $s_2=0$ and gives
\[
 z(a_h)=\frac{(1-\eta)\lambda}{1-\lambda}\ne0,
 \qquad z(b_h)=s_1=0,
\]
with every other handle coordinate zero.  The winding $a_h$
functional is therefore nonzero on this direction, with either
orientation of $a_h$.  The fixed-hyperplane description and Hermitian
symmetry give the reciprocal nonzero pairing, so the $a_h$-winding
direction has $s_1\ne0$.  That direction has zero coordinates on all
$a_j,b_j$ for $j<h$, by the last assertion of
Lemma~\ref{p4:winding-arcs}; subtracting its gauge coboundary does not
alter these coordinates, whose multipliers are one.  Hence every
remaining winding functional is nonzero on this direction.  The graph
is connected.  All these pair twists have exceptional eigenvalue
$\lambda\lambda^{-1}=1$, so their nontrivial changes are rank-one
nilpotents.  Lemmas~\ref{p4:twist-model} and
\ref{p4:connected-irreducibility} finish this case.

In all cases the connected twist closure is irreducible and its Lie
algebra contains a rank-one nilpotent.  It lies in the projected Hodge
group at every $s\in\mathcal T_\phi^{\mathrm{Hg}}$ by
Lemma~\ref{p4:generic-tensors}.  The latter is
reductive and irreducible, so Lemma~\ref{p4:cover-rank-one} applied to it
gives the stated standard special linear or symplectic image.

The same constructions apply after every scalar conjugation: the
active set is unchanged, a multiplier equal to one remains one, and
every displayed nonzero factor remains nonzero.  Thus every conjugate
space has a nontrivial unipotent preserving its polarization Hermitian
form.  A definite Hermitian space admits no such unipotent isometry.
Polarization is definite on a pure character space, so both Hodge
types occur on every conjugate.

For the trivial character use the ordinary handle transvections on
$H^1(B)$, again taking powers in the finite-index stabilizer.  They
give the full symplectic group by the same spanning and connectedness
argument and Lemma~\ref{p4:cover-rank-one}; the polarization rules out the
larger special linear alternative in dimension greater than two.
\end{proof}


\subsection{Characters over the sphere}
\label{p4:sphere}

In this subsection the base has genus zero.  Let $P$ be the finite
abelian deck group, let $\chi:P\to k^\times$, $k=\Qbar$, be a
nontrivial character, and put
\[
 W_\chi=H^1(C,k)_\chi.
\]
As in Lemma~\ref{p4:twist-model}, write $t$ for the corresponding
rank-one local system, discard its inactive punctures when computing
cohomology, and denote its active local multipliers by
$t_1,\ldots,t_r$.  Thus
\begin{equation}\label{p4:sphere-dimension}
 t_i\ne1,\qquad \prod_{i=1}^r t_i=1,\qquad
 n:=\dim_k W_\chi=r-2.
\end{equation}
Additional marked points can be retained in the parameter space.  They
do not alter $W_\chi$.  Let
$M_\chi\subset\operatorname{GL}(W_\chi)$ be the Zariski closure of the
already defined $\Gamma$-action.
All the two-puncture twists used below belong to this stabilizer:
their substitutions conjugate local meridians, which leaves the
homomorphism to the abelian group $P$ unchanged.  Changing the local
lift used to compute a formula multiplies its character action by a
root of unity and has no effect on the projective comparisons or Lie
algebras.

Fix $s\in\mathcal T_\phi^{\mathrm{Hg}}$ and write $J=J_s$.
A character space is
\emph{non-toric} if the image of $J^{\mathrm{der}}$ on it is nontrivial.
Statements about all scalar translates refer to all embeddings into
$\C$ of a number field over which the indicated spaces and maps are
defined, or, equivalently, to all their extensions to $k$.

For rank at least three, finite monodromy forces the Hodge action on the
rational scalar orbit to be toric, while infinite monodromy yields a
standard derived projection.  Ranks at most two are handled directly at the
Hodge-group level.  Only after these projections have been determined
do we compare character spaces on which the same simple ideal acts.

\begin{lemma}[Finite determinant image and scalar orbits]\label{p4:sphere-determinant-finite}
The determinant character of $M_\chi$ has finite image.  If the
monodromy on one scalar character space is finite, the monodromy on
the rational sum of its scalar orbit is finite as well.
If the Hermitian forms on every member of a scalar orbit are definite,
the monodromy on that rational sum is finite.
\end{lemma}

\begin{proof}
Choose a cyclotomic number field $E$ containing all character values.
The projector for $\chi$ and the monodromy representation on its image
are defined over $E$.  The determinant of an integral monodromy
operator on this image is an algebraic integer: its eigenvalues are
among those of the operator on integral first cohomology.  Applying
the same observation to the inverse shows that the determinant is an
algebraic unit.  At every embedding of $E$ the operator preserves the
nondegenerate Hermitian form obtained from the polarization, so the
absolute value of its determinant is one.  An algebraic integer all
of whose conjugates have absolute value one is a root of unity.
For example, the monic polynomials whose roots are the conjugates of
its positive powers have bounded integral coefficients; their finite
number forces two powers of a conjugate, and hence of the original
integer, to coincide.  The roots of unity in the fixed field $E$ form
a finite group.  This finite root set is Zariski closed, so the
determinant of the closure $M_\chi$ has the same finite-image property.
This proves the first assertion.

The matrices at any other embedding are obtained by applying that
embedding to the matrices over $E$.  A finite set of matrices therefore
remains finite after every embedding.  Their direct sum is the scalar
extension of the rational orbit representation, proving the second
assertion.  For the last assertion, change the sign of each negative
definite Hermitian form and take their direct sum.  It gives an
invariant positive definite real form on the rational orbit
representation, after restricting from its complexification.  This
representation also has an invariant lattice, its intersection with
integral first cohomology.  The image of each lattice basis vector has
fixed length and hence finitely many possibilities.  The monodromy
image is finite.
\end{proof}

The finite case can now be separated from the standard group actions.
Its decisive feature is that the full branch variation rules out a
mixed character with constant period.

\begin{lemma}[Finite monodromy and purity]\label{p4:sphere-finite-pure}
Suppose $n\ge2$.  If $M_\chi$ is finite, every scalar translate of
$W_\chi$ is pure at every fiber.  At every fiber its rational Hodge
group acts through scalar operators on each character space in that
orbit, and the action on their rational sum is toric.
\end{lemma}

\begin{proof}
Let $V_\Q$ be the rational character-orbit summand.  The orbit contains
the inverse character, so the polarization restricts nondegenerately
to $V_\Q$.  It is a polarized subvariation of first cohomology.
By Lemma~\ref{p4:sphere-determinant-finite} its monodromy is finite.
We pass to an algebraic quotient of the full marked family.
Since $r\ge4$, the ordered moduli space of the marked points on
$\mathbb P^1$ is a smooth quasi-projective configuration space modulo
projective coordinates; one may fix three labels.  Include the labelled
base point and the chosen point above it.  The finite-index stabilizer
of $\phi$ gives a finite topological cover of this configuration
space, which is algebraic and finite \'etale by the Riemann existence
theorem.  The base-point section splits the fundamental-group sequence
of the universal punctured curve: in the homotopy sequence of this
surface bundle, the section kills the boundary map into the fiber
fundamental group and splits the projection on fundamental groups.
Since the base subgroup preserves
$\phi$, sending its loops to the identity and the fiber loops by
$\phi$ defines a homomorphism from the total fundamental group to $P$.
Its finite topological cover is algebraic by the same theorem.
Normalize the universal compact curve in this cover.  Analytically
near a branch section, a small simply connected base neighborhood
trivializes the topological data, and the normalization is a disjoint
union of charts $(u,s)\mapsto(u^e,s)$ with the fixed inertia order.
The normalization is finite over the universal compact curve, so these
charts make it a smooth proper family of compact covers.  Thus this
is a quotient of the full marked variation, not a restricted family.
Indeed its monodromy on first cohomology is the action of the chosen
lifts in $\Gamma$: the lifted base-point section fixes the sheet during
transport.  On pulling back to $\mathcal T_\phi$, it is the original
marked family, since the topological covering epimorphism and the
chosen sheet agree and the complex structure is uniquely obtained by
lifting the base charts and filling the same power charts.  In
particular the pulled-back orbit variation is exactly $V_\Q$ above.

A finite-index kernel in $\Gamma$ kills the finite monodromy of
$V_\Q$.  The corresponding finite topological cover of the same
configuration quotient is algebraic and finite \'etale by Riemann
existence; the smooth proper family is simply pulled back to it.
The period map of $V_\Q$, now with a constant marking,
is a single-valued holomorphic map from this cover to a Siegel
period domain.  A Cayley transform realizes the domain as a bounded
domain in a finite-dimensional complex vector space.

Each coordinate of the period map is consequently a bounded
holomorphic function.  Take a smooth projective compactification
with normal-crossing boundary.  The removable-singularity theorem,
in local polydisks along the boundary, extends each coordinate
holomorphically over the compactification.  Such a function is
constant on a connected projective variety.  The period map of
$V_\Q$ is therefore constant.

We show that a mixed character would contradict this constancy.
Pass to the cyclic quotient
\[
 \pi:Y=C/\ker\chi\longrightarrow\mathbb P^1
\]
and let $D_\chi$ consist of its $r$ active branch points.  The $\chi$ and
$\chi^{-1}$ spaces in first cohomology are unchanged by this
quotient.  The map forgetting inactive and auxiliary markings is
locally a submersion onto the configuration space of $D_\chi$: the
active points can be moved independently while the additional
points are kept away from them.  It thus suffices to show that the
character period on the cyclic quotient has nonzero differential
whenever it is mixed.

Here is a deformation calculation which also covers genus one for
$Y$.  If a branch chart is $x=u^e$, an invariant vector field on
$Y$ has the form $u b(u^e)\partial/\partial u$; its pushforward is
$e x b(x)\partial/\partial x$.  Consequently
\begin{equation}\label{p4:sphere-tangent}
 (\pi_*T_Y)^{P/\ker\chi}=T_{\mathbb P^1}(-D_\chi),\qquad
 H^1(Y,T_Y)^{P/\ker\chi}
   =H^1(\mathbb P^1,T_{\mathbb P^1}(-D_\chi)).
\end{equation}
The second equality uses finiteness of $\pi$ and exactness of
invariants in characteristic zero.  The tangent map for varying the
branch points is this identification.  Indeed, on a \v{C}ech cover,
coordinate changes preserving the branch sections lift to the
cyclic cover; differentiation gives precisely the inverse of the
displayed pushforward on invariant vector fields.  This describes
the full tangent space of ordered branch configurations modulo
projective coordinate changes.  Since $r\ge4$, there are no
infinitesimal automorphisms of the pointed base.  This argument does
not require $H^0(Y,T_Y)=0$.

If the $\chi$-space is mixed, there are nonzero forms
\[
 \alpha\in H^0(Y,K_Y)_\chi,
 \qquad \beta\in H^0(Y,K_Y)_{\chi^{-1}}.
\]
Their product $\alpha\beta$ is a nonzero invariant quadratic
differential.  It pairs nontrivially with some element of
$H^1(Y,T_Y)^{P/\ker\chi}$ by Serre duality and averaging over the
finite group.  Equivalently, invariant quadratic differentials
descend to $H^0(\mathbb P^1,K_{\mathbb P^1}^{\otimes2}(D_\chi))$:
in the chart $x=u^e$, invariance permits at most a simple pole at
$x=0$, and its pullback is holomorphic.  This is the dual of
\eqref{p4:sphere-tangent}.

For clarity, the period differential here follows directly from
the same local deformation.  If the transition functions of the
deformed curve differ to first order by the vector-field cocycle
$\theta_{ij}$, Cartan's formula for the closed holomorphic form
$\alpha$ shows that its first-order changes differ on overlaps by
$\mathcal L_{\theta_{ij}}\alpha
 =d(\theta_{ij}\mathbin{\lrcorner}\alpha)$.
Its variation modulo holomorphic forms is therefore the class of
$\theta_{ij}\mathbin{\lrcorner}\alpha$ in $H^1(Y,\mathcal O_Y)$,
up to an immaterial common sign.  Serre duality pairs this class with
$\beta$ by evaluating $\theta$ on $\alpha\beta$.  The chosen tangent
direction makes that value nonzero.  This calculation also applies
in genus one.  Thus the period differential is nonzero.  The
contradiction proves purity, separately at every scalar translate.

Finally write $F=\Q(\chi(P))$.  The rational orbit summand is an
$F$-vector space whose scalar embedding spaces are exactly these
character spaces.  Since each is pure, the Hodge circle acts by
scalars on each of them.  It therefore centralizes the entire
rational algebra $\operatorname{End}_F(V_\Q)$.  The definition of the rational
Hodge group at each fiber forces it to centralize the same algebra.  Its action
on $V_\Q$ lies in the scalar torus
$\operatorname{Res}_{F/\Q}\mathbb G_m$, as claimed.
\end{proof}

It remains to identify the infinite character monodromy.  We first
prove connected irreducibility and produce the rank-one Lie element
needed by Lemma~\ref{p4:cover-rank-one}.

\begin{lemma}[Infinite sphere monodromy and rank-one Lie elements]\label{p4:sphere-connected-irreducible}
Suppose $n\ge3$.  If $M_\chi$ is infinite, its identity component
acts irreducibly on $W_\chi$.
Moreover its Lie algebra contains a rank-one nilpotent.
\end{lemma}

\begin{proof}
We first check that the arc family acts irreducibly.  When an actual
rank-one unipotent is available, its logarithm yields the nilpotent needed
to prove connected irreducibility; the remaining pseudoreflection case
requires a separate proof of connected irreducibility.  In that case an
adjoint-eigenspace projection then produces a rank-one nilpotent in the
Lie algebra.

Consecutive two-puncture arcs give an irreducible family on
$W_\chi$.  To check the hypotheses of
Lemma~\ref{p4:twist-model}, use the cocycle coordinates
$z_i=z(g_i)$ and put $s_i=z_i/(1-t_i)$.  Their common fixed
hyperplanes say that all $s_i$ are equal; subtracting the resulting
coboundary kills the cocycle.  Consecutive arc directions have a
nonzero mutual pairing by the formula in that lemma.  Thus the
directions span and their nonorthogonality graph is connected.
The full group $M_\chi$ acts irreducibly.

Its identity component is reductive in this representation.
Indeed, the unipotent radical is characteristic in the identity
component and hence normal in $M_\chi$.  Its nonzero fixed-vector
space is then $M_\chi$-stable, and irreducibility makes that space
all of $W_\chi$.  The radical acts trivially.  The isotypical
summands for $M_\chi^\circ$ are permuted transitively by
$M_\chi$.

First suppose that the Lie algebra already contains a rank-one
nilpotent $N$.  On an isotypical summand $S\otimes A$ it acts as
$N_S\otimes1_A$.  A summand on which it is nonzero must therefore
have $\dim A=1$ and $\dim S\ge2$.  Transitivity gives these same
properties on every isotypical summand.  An operator differing
from a scalar by rank one cannot permute two summands of dimension
at least two: its displacement on either of the summands would
have rank at least two.  Every arc pseudoreflection therefore fixes
each summand.  The irreducibility of the arc family forces there
to be just one, proving connected irreducibility in this case.

There are two immediate sources of such an $N$.  If $t_it_j=1$
for some active pair, its twist is a nontrivial rank-one unipotent
by Lemma~\ref{p4:twist-model}.  Its logarithm belongs to
$\operatorname{Lie}(M_\chi^\circ)$.  If all active pair products are $-1$, then
all the $t_i$ are equal to the same primitive fourth root of unity.
Equation~\eqref{p4:sphere-dimension} implies $4\mid r$; since
$r\ge5$, we have $r\ge8$.  Choose a disk containing four active
points.  If $g$ is the product of their consecutive meridians,
then $t(g)=1$.  The boundary twist conjugates each inside meridian
by $g$, and its cocycle displacement is
\[
 z_i\longmapsto z_i+(1-t_i)z(g)\quad(i\text{ inside}),
 \qquad z_j\longmapsto z_j\quad(j\text{ outside}).
\]
This is a rank-one unipotent.  Its image is nonzero modulo
coboundaries because there are active points outside the disk.
Its functional is nonzero: one can prescribe a nonzero value of
$z(g)$ on the inside and satisfy the global cocycle relation using
an outside entry.  Its square-zero displacement follows also from
$t(g)=1$.

It remains to treat the case with no pair product equal to $1$
and with some pair product $\lambda\ne\pm1$.  All pair twists
are then semisimple pseudoreflections.  Suppose first that the
isotypical summands for $M_\chi^\circ$ have dimension at least
two.  The rank argument just given makes every pair twist fix
every summand, so there is just one summand $S\otimes A$.
If $\dim S=1$, the connected group acts by scalars; its finite
determinant then makes the connected group trivial, contrary to
infinitude.  Thus $\dim S>1$.

An operator normalizing this isotypical action has the form
$B\otimes D$, with $B\in\operatorname{GL}(S)$ and $D\in\operatorname{GL}(A)$.
To see this, choose an intertwiner $B$ for the induced automorphism
of the irreducible representation on $S$; what remains centralizes
that representation and is $1\otimes D$ by Schur's lemma.
A nontrivial semisimple pseudoreflection cannot have this form
when both factors have dimension at least two.  Indeed the factors
must themselves be semisimple (a nontrivial unipotent Jordan part
would survive on the tensor product).  Their eigenvalue products
would then form a rectangular grid with just one exceptional
entry.  A two-by-two subrectangle containing that entry contradicts
the equality of the products of its two diagonal pairs.  Hence
$\dim A=1$, and the connected action is irreducible.

The only remaining reducible case has distinct one-dimensional
isotypical summands.  The Hermitian form pairs each such line
nondegenerately with exactly one other line: for characters
$a,b$ of the connected group, a nonzero pairing requires
$a\overline b=1$ on the finite-index monodromy subgroup dense in
that group.  Distinct characters preclude two choices of $b$.
Hermitian symmetry makes this a pairing involution on the set of
lines.  Each pair twist preserves the involution and permutes the
lines by either the identity or one transposition.  Indeed, a
permutation cycle of length $\ell$ contributes rank at least $\ell-1$
to the difference from a scalar identity, and disjoint cycles
contribute independently; here that rank is one.  The permutation
group generated by the twists is transitive, by irreducibility.
If the pairing involution had no fixed point, a transposition
commuting with it would have an invariant two-element support, so
could only exchange the members of one of its pairs.  Such
transpositions cannot act transitively on at least
three lines.  Thus there is a fixed point, and transitivity makes
every line self-paired.  Distinct lines are orthogonal; transitivity
also makes their signs equal.  The Hermitian form is definite.

The decomposition into distinct one-dimensional connected
isotypicals and the pseudoreflection action are algebraic
properties over $k$.  They persist at every scalar translate.
The same argument therefore makes the Hermitian form definite on
every translated character space.  By
Lemma~\ref{p4:sphere-determinant-finite} the rational orbit
monodromy is finite, a contradiction.  This proves connected
irreducibility in the remaining case.

Use now a pair twist with eigenvalue ratio
$\lambda\ne\pm1$.  After a scalar normalization it is
$R=\operatorname{diag}(\lambda,1,\ldots,1)$ for a decomposition
$W_\chi=L\oplus U$, $\dim L=1$.  Conjugation by $R$ preserves
$\mathfrak m=\operatorname{Lie}(M_\chi^\circ)$ and has the three distinct
eigenvalues $1,\lambda,\lambda^{-1}$ on
\[
 \operatorname{End}(L)\oplus\operatorname{End}(U),\qquad
 \operatorname{Hom}(U,L),\qquad \operatorname{Hom}(L,U),
\]
respectively.  Polynomial projection onto these eigenspaces
preserves $\mathfrak m$.  Connected irreducibility implies that
some off-diagonal component is nonzero, since otherwise $L$ and
$U$ would be invariant.  A nonzero element of either off-diagonal
space has rank one and square zero.  It supplies the asserted
nilpotent.
\end{proof}

The preceding lemma determines the infinite projections.  The finite
case is toric by Lemma~\ref{p4:sphere-finite-pure}.
For cyclic covers of $\PP^1$, the monodromy projections for mixed
characters with $\chi^2\ne1$ are computed in
\cite[preprint, Theorem~5.1.1]{Rohde2009};
\cite[Proposition~4.7]{Moonen2010} also includes the symplectic
projection for order-two characters.  The following
statement collects exactly the alternatives needed for the tensor
criterion.

\begin{proposition}[Sphere projections]\label{p4:sphere-projections}
At every $s\in\mathcal T_\phi^{\mathrm{Hg}}$ the following assertions hold.
\begin{enumerate}[label=\textup{(\roman*)}]
\item If $n\ge3$, then either $M_\chi$ is finite and the Hodge
action on its rational scalar orbit is toric, or
\[
 M_\chi^\circ=
 \begin{cases}
  \operatorname{Sp}(W_\chi),&\chi^2=1,\\
  \operatorname{SL}(W_\chi),&\chi^2\ne1.
 \end{cases}
\]
In the infinite case the derived Hodge projection is the same
standard group.
\item In dimensions at most two every non-toric character space
has dimension two and has standard derived Hodge projection
$\operatorname{SL}_2$.
\end{enumerate}
Here the symplectic group is for the restricted polarization on
the order-two character space.  Each non-toric character space
is supported by a single simple ideal of $\operatorname{Lie}(J^{\mathrm{der}})_k$.
\end{proposition}

\begin{proof}
For $n\ge3$, Lemmas~\ref{p4:sphere-connected-irreducible} and
\ref{p4:cover-rank-one} give either the full standard special linear or
the full standard symplectic derived algebra of connected
monodromy.  Its determinant is finite, so its connected scalar
center is trivial.  In the symplectic case all of $M_\chi$
normalizes the connected symplectic group and hence preserves
its unique invariant alternating form up to scalar.

In dimension at least four a semisimple pseudoreflection with
ratio different from $1$ cannot be such a similitude.  Normalize
its fixed hyperplane $U$ to have eigenvalue $1$.  The alternating
form does not vanish identically on $U$, since $\dim U=n-1>n/2$;
thus the similitude multiplier is $1$.  Its exceptional eigenline
must then be orthogonal to $U$ and to itself, and hence be in the
radical, an impossibility.  It follows that every active pair
product is $1$.  With at least three active points this forces
all $t_i=-1$.  The sphere meridians generate the deck group, so
$\chi$ has order two.  Conversely an order-two character space
has a nondegenerate alternating polarization and monodromy
contained in its symplectic group.  Its pair twists are nontrivial
unipotents, so the symplectic alternative occurs.  This proves
the stated connected monodromy groups.

The algebraic image of $L_k$ on $W_\chi$ is closed and is the Zariski
closure of the character action, hence is $M_\chi$.  The image of
$L_k^\circ$ is $M_\chi^\circ$, by the identity-component property of
algebraic group images.  By Lemma~\ref{p4:generic-tensors},
$L^\circ\subset J$, so connected character monodromy is contained in
the image of $J_k$ on $W_\chi$.  In the special linear case its action already contains
the full traceless algebra.  In the order-two case the Hodge group
also preserves the restricted alternating polarization.  Hence
the derived Hodge projections are exactly the groups asserted
in (i).  Finite monodromy gives the toric conclusion by
Lemma~\ref{p4:sphere-finite-pure}.

A connected reductive subgroup of $\operatorname{GL}_2$ either acts through a
torus or, when its derived image is nontrivial, has the full
standard $\operatorname{SL}_2$ as its derived image.  Indeed a reducible
two-dimensional representation of a reductive group is a sum of
characters, while the only nonzero semisimple subalgebra of
$\mathfrak{sl}_2$ is $\mathfrak{sl}_2$ itself.  Dimension one is
toric, and dimension zero is omitted.  This proves (ii).
Finally each displayed derived image is simple, and a
surjective map from a semisimple Lie algebra to a simple one is
supported by exactly one simple ideal.
\end{proof}

We have identified each projection, but a single simple Hodge ideal
may act on several character spaces.  To identify those copies by
algebraic maps, we must compare the whole mapping-class actions.
The only additional numerical calculation occurs for four active
points, where the standard spaces have dimension two.

\begin{lemma}[Four-point signatures]\label{p4:sphere-four-signatures}
Suppose $r=4$ and write $t_i=\exp(2\pi i v_i)$ with
$0<v_i<1$.  Set $S=\sum_i v_i\in\{1,2,3\}$.  In a fixed
choice of character convention the two Hodge multiplicities are
$S-1$ and $3-S$.  Thus $S=2$ means mixed, while $S=1$ and
$S=3$ give the two opposite pure types.
\end{lemma}

\begin{proof}
Pass to the cyclic quotient on which the character is faithful,
of order $m$, and write $v_i=a_i/m$, $0<a_i<m$.
Choose a coordinate with infinity away from the branch points.
The quotient has an equation
\[
 y^m=\prod_i(x-b_i)^{a_i},\qquad \sum_i a_i=mS.
\]
The eigendifferentials for the inverse of the character of $y$
are $f(x)\,dx/y$.  At a branch point with
$d_i=\gcd(m,a_i)$, a uniformizer has
$x-b_i=u^{m/d_i}$ and $y$ equal to a unit times
$u^{a_i/d_i}$.  The differential $dx/y$ is holomorphic there,
whereas a pole of $f$ would destroy holomorphicity.  Thus $f$ is
a polynomial.  At infinity holomorphicity requires
$\deg f\le S-2$, giving dimension $S-1$.  Replacing $y$ by
$\prod_i(x-b_i)/y$ computes the opposite eigenspace with
exponents $m-a_i$ and dimension $3-S$.  The choice of whether
$t$ labels this character or its inverse exchanges the two
dimensions everywhere and does not affect the assertion.
\end{proof}

\begin{proposition}[Links within one sphere cover]\label{p4:sphere-links}
Suppose two non-toric character spaces $W_\chi,W_{\chi'}$
are supported by the same simple Hodge Lie ideal.
They have the same dimension $n$ and standard derived
representations equivalent up to duality.
\begin{enumerate}[label=\textup{(\roman*)}]
\item If $n\ge3$, then $\chi'=\chi$ or $\chi^{-1}$.
The inverse character is identified by the polarization with the
polarized dual.  If the standard group is symplectic, the
character has order two and the space is mixed at every
scalar translate.
\item If $n=2$, one can choose either $W_\chi$ or its polarized
dual as model so that an equivalence of its derived
representation with $W_{\chi'}$ respects the Hodge circle at
every scalar translate.
\end{enumerate}
\end{proposition}

\begin{proof}
Write $\mathfrak f$ for the shared ideal.  The standard
representations of a simple special linear or symplectic
algebra determine their dimension and are preserved, up to
duality, by isomorphisms of these algebras.  The coincidence
$\mathfrak{sl}_2=\mathfrak{sp}_2$ has the same standard
two-dimensional representation.  Invert one character when
necessary, so that an $\mathfrak f$-linear isomorphism
\[
 F:W_\chi\longrightarrow W_{\chi'}
\]
exists.  At this stage $F$ need not respect the Hodge circle.

The projective actions of the entire group $\Gamma$ are
intertwined by $F$.  Indeed, $\Gamma$ preserves the character
spaces and normalizes $J$ by Lemma~\ref{p4:generic-tensors}.
The kernel of either character projection on
$\operatorname{Lie}(J^{\mathrm{der}})_k$ is preserved, because the
character space is preserved.  Its unique complementary supporting
ideal $\mathfrak f$ is therefore preserved as well.  Each lift induces
the same automorphism of this ideal on both spaces.  The two
implementations of this automorphism differ, after conjugation
by $F$, by a scalar, by irreducibility and Schur's lemma.

A two-puncture twist is projectively trivial if at least one
endpoint is inactive: fill that inactive puncture in the
rank-one local system, after which the twist is supported in
a disk with at most one puncture.  It is projectively nontrivial
when both endpoints are active.  For the latter assertion use
$r=n+2\ge4$ and the nonzero rank-one formula of
Lemma~\ref{p4:twist-model}; the global cocycle relation can be
satisfied using an active point outside the disk, so its
functional is nonzero as well as its direction.  The graph of
nontrivial pair twists is therefore the complete graph on the
active set.  Projective equivalence shows that the active sets
of $\chi$ and $\chi'$ coincide.  Label them identically and
write $t_i,t_i'$ for their multipliers.

Suppose first that $n\ge3$.  The pair twist has one exceptional
eigenvalue and a fixed hyperplane.  Its ratio is distinguished
projectively by the multiplicities, so
\begin{equation}\label{p4:sphere-pair-products}
 t_i't_j'=t_it_j\qquad(i\ne j).
\end{equation}
If the ratio is $1$, the operator is nontrivial unipotent and
the same conclusion holds.  These twists also intrinsically
identify their image directions: for a semisimple twist this
is its exceptional eigenline, and for a unipotent twist it is
the image of its normalized displacement.

Choose three active points in an embedded disk, with consecutive
meridians $g_i,g_j,g_\ell$.  Cocycles supported on these entries
and with value zero on their product form a two-dimensional
subspace $U\subset W_\chi$.  There are active points outside
the disk, so passage to the cocycle quotient is injective on
this space.  The two consecutive pair directions are independent
and span $U$.  Thus $F$ identifies this space with its primed
counterpart.  The twist enclosing the three points acts on
$U$ by the scalar $t_it_jt_\ell$ and on $W_\chi/U$ by $1$.
Indeed its displacement on a cocycle is supported in the disk
and has value zero on the product; on a cocycle with that
product value zero, conjugation by the product multiplies
every inside value by its multiplier.  The quotient has
dimension $n-2>0$.  Its identity action fixes the projective scalar
in the comparison, even when the full operator is a nontrivial
unipotent extension of the identity.  We obtain
\[
 t_i't_j't_\ell'=t_it_jt_\ell.
\]
Equation~\eqref{p4:sphere-pair-products} makes every ratio
$t_i'/t_i$ equal to a common sign; the triple equality makes
that sign $1$.  All local multipliers are equal.  Since sphere
meridians generate the deck group, the characters coincide
after the permitted inversion.  This proves the first part
of (i).  The symplectic assertion follows from
Proposition~\ref{p4:sphere-projections}: its character has
order two, so complex conjugation acts internally and exchanges
its two Hodge parts.  Both parts have dimension $n/2$ at every
translate.

Now suppose $n=2$, so the common active set has four points.
The eigenvalue ratio of a projective two-by-two operator is
defined up to inversion.  Hence
\begin{equation}\label{p4:sphere-four-pair-products}
 t_i't_j'\in\{t_it_j,(t_it_j)^{-1}\}.
\end{equation}
Put $A=t_1t_2$, $B=t_1t_3$, and $D=t_1t_4$.
The Klein four permutations on the four indices give the
following triples of pair products:
\[
\begin{array}{c|c}
 \pi&(t_{\pi(1)}t_{\pi(2)},
       t_{\pi(1)}t_{\pi(3)},t_{\pi(1)}t_{\pi(4)})\\ \hline
 1&(A,B,D)\\
 (12)(34)&(A,B^{-1},D^{-1})\\
 (13)(24)&(A^{-1},B,D^{-1})\\
 (14)(23)&(A^{-1},B^{-1},D)
\end{array}
\]
Uniformly inverting the four entries supplies the other four
inversion patterns.  Choose $\pi$ in this Klein four group and
$\epsilon\in\{1,-1\}$ realizing the three choices in
\eqref{p4:sphere-four-pair-products} involving index $1$.
Since $t_1^2=ABD$, recovery of the four entries from these
three products leaves only a common sign.  Thus
\begin{equation}\label{p4:sphere-four-tuple}
 t_i'=\delta\,t_{\pi(i)}^\epsilon\quad(1\le i\le4),
 \qquad \delta\in\{1,-1\}.
\end{equation}
The choices $\delta,\epsilon,\pi$ are made once and are
unchanged at every scalar translate of this identity.

This classifies only the projective twist data.  It remains to use the
shared ideal to align pure and mixed behavior and to make one fixed choice
between the original model and its polarized dual for every scalar translate.

Mixedness agrees between the two character spaces at each
translate.  To justify this without any assumption about the
pure orientations, split the infinitesimal Hodge-circle action
into its central scalar part and its component in the shared
simple ideal.  On a standard two-dimensional module, the
ideal component is zero precisely when the circle is scalar,
that is, precisely when the character space is pure.  The
same ideal component acts on both equivalent modules, and
this observation remains valid after every scalar translation.

We compare the two pure orientations.  Permuting the four
multipliers has no effect on their signatures; uniformly
inverting them reverses the two pure signatures by
Lemma~\ref{p4:sphere-four-signatures}.  Consider therefore
the additional common sign $\delta=-1$.  At a fixed translate
write the four entries before this sign change as
$\exp(2\pi i v_i)$, $0<v_i<1$, and put
$S=\sum_i v_i$.  No $v_i$ equals $1/2$, since otherwise
the sign change would produce an inactive multiplier.  If
$a$ entries exceed $1/2$, the sum after the change is
\begin{equation}\label{p4:sphere-sign-change}
 S'=S+2-a.
\end{equation}
For $S=1$ one has $a\in\{0,1\}$; agreement of purity
excludes $S'=2$, so $a=0$ and $S'=3$.
For $S=3$ one has $a\in\{3,4\}$; the same argument gives
$a=4$ and $S'=1$.  For $S=2$, agreement of mixedness forces
$a=2$ and $S'=2$.  Thus the common sign reverses pure
orientation wherever the spaces are pure.  Combining it
with $\epsilon$, the relative orientation of the two pure
spaces is the same fixed sign $\delta\epsilon$ at every
pure scalar translate.

The standard $\mathfrak{sl}_2$ module is self-dual.  According
to this fixed sign choose either the original model or its
polarized dual, whose circle weights are the negatives of
the original weights.  Its pure type now agrees with that
of $W_{\chi'}$ at every translate.  Choose an intertwiner
over $k$ of the resulting standard ideal modules.  At a
mixed translate the two circle weights are $+1,-1$; their
sum is zero, so the central scalar shift is zero on both
modules.  The intertwiner respects the common ideal action
and hence the whole circle.  At a pure translate both
circles are the same scalar, so it respects them there as
well.  The same reasoning applies to every scalar conjugate
of the chosen intertwiner.  This proves (ii).  Notice that
polarized duality reverses the weights at every translate;
no commutation between an arbitrary scalar automorphism and
complex conjugation is being assumed.
\end{proof}

\subsection{Linked character spaces and saturation}

The sphere comparison is now complete.  In positive genus every
nonzero character is mixed at all scalar conjugates, so the analogous
comparison follows from a shorter Hodge-circle argument.

\begin{lemma}[Links between mixed standard spaces]\label{p4:mixed-links}
Let two admissible irreducible spaces in the same polarized weight-one
structure have equivalent representations of a single supporting simple
ideal of the derived Hodge Lie algebra.  Suppose both spaces are mixed
at every simultaneous scalar conjugate.  An intertwiner between them
is an admissible Hodge isomorphism.  The analogous dual comparison is
obtained by replacing one space by its opposite paired space.

If the derived action on one such space is standard symplectic, its
invariant alternating tensor is a pair tensor on that space itself.
\end{lemma}

\begin{proof}
The intertwiner line is one-dimensional by Schur's lemma and is stable
under the Hodge group, since its derived group is normal.  If the Hodge
circle acts on this line by $z^a$, the intertwiner shifts the spectrum
of one space by $a$.  Both spectra are $\{-1,1\}$, so
$\{-1+a,1+a\}=\{-1,1\}$ forces $a=0$.  Scalar conjugation gives
the identical argument for the conjugated intertwiner and spaces.
It is therefore an admissible Hodge isomorphism, and hence algebraic
in the sense of Companion results~\ref{comp:tools}\textup{(ii)}.

An invariant nondegenerate alternating form also spans a Hodge-stable
line.  Its scalar weight must be zero: nondegeneracy pairs the spectrum
$\{-1,1\}$ with its reflection about half that weight, so the midpoint
is zero.  The inverse tensor consequently has type $(1,1)$, and so do
all its conjugates.  It is the required pair tensor.
\end{proof}

\begin{proof}[Proof of Proposition~\ref{p4:cover-saturation} and Theorem~\ref{p4:abelian-covers}]
Fix $s\in\mathcal T_\phi^{\mathrm{Hg}}$ and put $J=J_s$.
Set $A=\operatorname{Jac}(C_s)$ and choose an Abel map $a:C_s\to A$.
Its pullback $a^*:H^1(A,\Q)\to H_\Q$ is an isomorphism of rational
Hodge structures.  Its inverse is an algebraic degree-one operation
by Companion results~\ref{comp:tools}\textup{(ii)}.  Transport the admissible
character projectors $e_\chi$ through this isomorphism.  Their images are then
actual summands of $H^1(A,k)$, as required by the tensor criterion;
we use the same notation for the identified curve and Jacobian spaces.

Collect the rational Galois orbits of character spaces on which
$J^{\mathrm{der}}$ acts trivially into a rational Hodge summand $H_0$.
Its Hodge group is a torus: its projections on the orbits are tori,
and the group on their direct sum is a subgroup of their product.
Companion results~\ref{comp:tools}\textup{(iv)}
realize $H_0$ as the first cohomology of a CM isogeny factor.
For $h\ge1$, every remaining nonzero character space has standard
special linear or symplectic derived projection by
Proposition~\ref{p4:positive-projections}.  For $h=0$ the same
assertion, with the toric alternatives accounted for, is
Proposition~\ref{p4:sphere-projections}.

Each remaining character space is supported by exactly one simple
ideal of $\operatorname{Lie}(J^{\mathrm{der}})_k$.  Indeed its image is
a simple standard special linear or symplectic algebra.  In a
semisimple algebra mapping onto a simple one, exactly one simple ideal
has nonzero image; the commuting images of two nonzero ideals could
not both be that simple algebra.  The supporting ideal maps
isomorphically.  If two character spaces share it, their standard
representations are equivalent or dual.  The only coincidence
between the two group types is
$\mathfrak{sl}_2=\mathfrak{sp}_2$, with the same standard module.

In positive genus, Lemma~\ref{p4:mixed-links} makes these comparisons
admissible Hodge isomorphisms, using the opposite character for a dual
comparison.  On the sphere, Proposition~\ref{p4:sphere-links} gives
the same conclusion, including the simultaneous alignment of pure
scalar translates in dimension two.  Group the spaces with the same
supporting ideal into one block.  Treat a two-dimensional block as
linear.  A higher-dimensional symplectic block has its alternating
pair tensor: it comes from Lemma~\ref{p4:mixed-links} in positive
genus and from the polarization on an order-two character on the
sphere.  A linear block uses the polarization pairing of a character
and its opposite.  Since finite character values are roots of unity,
the opposite space is its coefficient-conjugate space
$\overline W$.

The simple-ideal decomposition gives simultaneous independent
standard actions on these blocks.  More explicitly, each ideal acts
on precisely the copies in its own block, with the chosen standard
or dual identifications, and acts trivially on the other blocks and
on $H_0$.  The connected image of $J^{\mathrm{der}}_k$ consequently
has the Lie algebra of the product of these standard groups.  The
connected groups have the same image, so this is the product required
in part \textnormal{(c)} of Companion results~\ref{comp:saturated}.
No independence of the central torus is needed.

A representative of every block is the transported image of an actual
character space on $C_s$.  Corollary~\ref{p4:abelian-determinants} supplies its determinant
singleton: a nontrivial character uses the faithful cyclic quotient
$C_s/\ker\chi$, and the invariant character uses the polarization.
The curve--Jacobian degree-one maps and all the block identifications
are algebraic operations by Companion results~\ref{comp:tools}\textup{(ii)}.
They transfer these determinant lines to the chosen Jacobian slots.
Thus every linear determinant singleton is CM-supplied; there are
no orthogonal blocks.  The data are saturated.

Companion results~\ref{comp:saturated} now prove the assertion for every
power after adjoining any CM $M$.  This is a pointwise application
after saturation, so $M$ does not enter the definition of
$\mathcal T_\phi^{\mathrm{Hg}}$.  Finally, the algebraic
curve--Jacobian correspondences identify their first cohomology, and
a point and the fundamental class supply the other two curve degrees.
External products of these maps and
Companion results~\ref{comp:transfer} transfer the tensor
conclusion to every $C_s^a\times M^b$.
\end{proof}

\section{Coordinate-root covers}\label{sec:roots}
This section determines character ranks from their support and
shows that projection to that support preserves the character
cohomology used by the Gale correspondence.

All varieties are over \(\C\), and character eigenspaces are taken
over \(\kk\).  We use the convention that \(\Q(1)\)
has weight \(-2\).

\begin{definition}\label{def:root}
Let \(m,n\geq2\) and \(0\leq d\leq m-2\).  A subspace
\(P\subset\C^m\) of dimension \(d+1\) is \emph{coordinate-general}
if every \(d+1\) coordinate functionals restrict to a basis of
\(P^*\).  Let \(\U_{d,m}\subset\Gr(d+1,m)\) be this open locus.
For \(P\in\U_{d,m}\), define
\[
 X_P=\{[x_i]\in\PP^{m-1}:(x_1^n,\ldots,x_m^n)\in P\},\qquad
 G_{m,n}=(\mu_n)^m/\mu_n,
\]
where the denominator is diagonal.  The phase action preserves \(X_P\),
and \(q_P:X_P\to\PP(P)\), \([x_i]\mapsto[x_i^n]\), is its
coordinate-power map.
\end{definition}

\begin{proposition}[Exact smooth locus]\label{prop:smooth-locus}
Let \(n\geq2\), \(1\leq c\leq m-1\), and let \(A\) be a rank-\(c\)
matrix with \(m\) columns.  Then \(X_A\) is a complete intersection
of dimension \(m-c-1\).  It is smooth if and only if every
\(c\)-column minor of \(A\) is nonzero.  This condition is equivalent
to coordinate generality of \(P=\ker A\).
\end{proposition}
\begin{proof}
The ring map \(\C[t_1,\ldots,t_m]\to\C[x_1,\ldots,x_m]\),
\(t_i\mapsto x_i^n\), is free with basis the monomials
\(\prod_i x_i^{e_i}\), \(0\leq e_i<n\).  It is flat, so the
independent linear equations for \(\ker A\) remain a regular
sequence after pullback.  The projective coordinate-power map is
finite and surjective.  Its inverse image of \(\PP(\ker A)\)
therefore has dimension \(m-c-1\), proving the first assertion.

At a point let \(I=\{i:x_i\ne0\}\).  The Jacobian has column
\(n a_i x_i^{n-1}\) for \(i\in I\) and zero otherwise.  If every
\(c\)-minor is nonzero, every at most \(c\) columns are independent.
The equation \(\sum_{i\in I}a_i x_i^n=0\) forces \(|I|>c\),
and the Jacobian has rank \(c\).  Conversely, a vanishing
\(c\)-minor gives a nonzero dependence on those columns.  Taking
\(n\)th roots of its nonzero coefficients and setting the other
coordinates to zero gives a point with Jacobian rank below \(c\).
Here \(n\geq2\) makes the zero coordinates give zero Jacobian
columns.  The Jacobian criterion proves the equivalence.

Projection \(\ker A\to\C^S\), for \(|S|=m-c\), is an isomorphism
exactly when no nonzero vector of \(\ker A\) is supported on the
complementary \(c\) coordinates.  This is the nonvanishing of that
complementary minor, proving the coordinate-generality assertion.
\end{proof}

\begin{lemma}[Cohomology outside the middle degree]\label{lem:ambient}
Let \(X_A\) be a smooth diagonal complete intersection as in
Proposition~\ref{prop:smooth-locus}, of dimension \(d\geq1\),
and put \(h=c_1(\mathcal O_{X_A}(1))\).  For \(0\leq k\leq2d\)
and \(k\ne d\),
\[
 H^k(X_A,\Q)=
 \begin{cases}
 \Q h^{k/2},&k\text{ even},\\
 0,&k\text{ odd}.
 \end{cases}
\]
In particular \(X_A\) is connected.
\end{lemma}
\begin{proof}
We first choose a smooth partial-intersection flag.  Let \(Q\) be the
row space of \(A\), of dimension \(c\).  For every \(j\leq c\)
and every \(j\)-element coordinate set \(I\), the projection
\(Q\to\C^I\) is surjective: a dependent set of \(j\) columns
could be extended to a dependent set of \(c\) columns, contrary
to Proposition~\ref{prop:smooth-locus}.  In the complete flag
variety of \(Q\), the condition that this projection restricts
to an isomorphism on the \(j\)-plane is a nonempty open condition.
It is nonempty by choosing a complement of the projection kernel
and extending it to a flag.  The flag variety is irreducible, so
the intersection of these finitely many nonempty opens is nonempty.
Choose such a flag \(0=Q_0\subset Q_1\subset\cdots\subset Q_c=Q\).

Let \(Y_j\) be the diagonal zero locus of \(Q_j\), with
\(Y_0=\PP^{m-1}\).  Every maximal minor of \(Q_j\) is nonzero,
so Proposition~\ref{prop:smooth-locus} makes \(Y_j\) smooth of
dimension \(m-j-1\).  Choose an equation basis adapted to the
flag.  Under the degree-\(n\) Veronese embedding of \(\PP^{m-1}\),
the next degree-\(n\) equation is a hyperplane, so \(Y_j\) is a
hyperplane section of \(Y_{j-1}\) in that embedding.

The Lefschetz hyperplane theorem in
\cite[Corollary~7.3, p.~41]{Milnor} says that for a projective
variety \(V\) of dimension \(k\) and a hyperplane containing its
singular locus, the inclusion of the section induces homology
isomorphisms in degrees below \(k-1\).  Here \(Y_{j-1}\) is
smooth, so its singular locus is empty.  Applying the theorem
successively gives
\[
 H_i(X_A,\Z)\xrightarrow{\sim}H_i(\PP^{m-1},\Z)\quad(i<d).
\]
Universal coefficients over \(\Q\) give the corresponding
cohomology restriction isomorphisms.  This gives connectedness
at \(i=0\) and the formula below \(d\).

Poincar\'{e} duality gives the same dimensions above \(d\).
In even degree the ambient class is nonzero: the complete
intersection has degree \(n^c\), so \(\int_{X_A}h^d=n^c>0\);
vanishing of any \(h^r\), \(0\leq r\leq d\), would force
\(h^d=0\).  Thus it spans the one-dimensional even cohomology
above the middle as well.
\end{proof}

\begin{lemma}[Root-cover geometry]\label{lem:geometry}
For \(P\in\U_{d,m}\), \(X_P\) is a smooth projective complete
intersection of dimension \(d\), connected if \(d\geq1\).
The map \(q_P\) is finite and is the quotient by \(G_{m,n}\).
The coordinate hyperplanes on \(\PP(P)\) and their reduced inverse
images on \(X_P\) have simple normal crossings.  For \(d=0\),
\(X_P\) is a \(G_{m,n}\)-torsor of reduced points.

The open \(\U_{d,m}\) is nonempty and connected, and the \(X_P\)
form a smooth projective family over it with its phase action.
\end{lemma}
\begin{proof}
Any at most \(d+1\) coordinate restrictions to \(P\) are independent.
Thus at most \(d\) coordinates vanish at a point of \(\PP(P)\).
If \(a\) vanish, their ratios to a nonvanishing coordinate have
independent differentials and extend to relative local coordinates
\(t_1,\ldots,t_d\).  After an \'{e}tale base change taking roots of
the remaining units, the cover has equations
\[
 z_1^n=t_1,\ldots,z_a^n=t_a
\]
and free coordinates \(t_{a+1},\ldots,t_d\).  It is smooth, with
normal-crossing boundary \(z_1\cdots z_a=0\).  The same calculation
on the universal projective bundle proves the relative assertions.
Proposition~\ref{prop:smooth-locus} gives the complete-intersection
description.

Fibers of the coordinate-power map are exactly phase orbits, giving
the finite quotient.  In dimension zero the base point has no zero
coordinate, so its inverse image is a free transitive phase set.
For \(d\geq1\), Lemma~\ref{lem:ambient} gives connectedness and
ambient cohomology outside the middle degree.

For distinct \(a_1,\ldots,a_m\in\C\), the subspace
\(\{(f(a_1),\ldots,f(a_m)):\deg f\leq d\}\) has nonzero maximal
minors by the Vandermonde determinant.  Thus \(\U_{d,m}\) is a
nonempty open in an irreducible Grassmannian, hence connected as a
complex manifold.  The universal incidence equations are projective
over it; the local calculation proves smoothness and respects phases.
\end{proof}

\subsection{Evaluation covers and character ranks}
Identify a character of \(G_{m,n}\) with
\[
 b=(b_1,\ldots,b_m)\in(\Z/n)^m,\quad \sum_i b_i=0,\quad
 \chi_b([\zeta_i])=\prod_i\zeta_i^{b_i}.
\]
Subscripts \(b\) mean eigenspaces for pullback \(g^*\).  Put
\(S(b)=\{i:b_i\ne0\}\), \(s(b)=|S(b)|\), and
\(W_{P,b}=H^d(X_P,\kk)_b\).  Full support means \(s(b)=m\).

In positive dimension we first use subspaces obtained by evaluating
polynomials at distinct points.  Multiplication of linear polynomials
realizes the corresponding root cover as a finite quotient of a power
of one curve.  The phase action and the
signs from permuting degree-one factors then identify its nontrivial
character cohomology with an exterior power.  This is the graded
alternating-character mechanism of
\cite[Lemmas~1.1--1.2]{Schoen}; here we identify the quotient explicitly.

\begin{lemma}[Evaluation quotient]\label{lem:evaluation}
Choose distinct \(a_1,\ldots,a_m\in\C\), \(m\geq3\), and put
\[
 P_j=\{(f(a_1),\ldots,f(a_m)):\deg f\leq j\}
 \quad(1\leq j\leq m-2),\qquad C=X_{P_1}.
\]
Coordinate multiplication defines a finite surjective morphism
\[
 \mu_j:C^j\to X_{P_j},\qquad
 (x^{(1)},\ldots,x^{(j)})\mapsto
       [\prod_{\nu=1}^j x_i^{(\nu)}]_{i=1}^m,
\]
which identifies the target with
\[
 C^j\big/\left(\ker(G_{m,n}^j\xrightarrow{\prod}G_{m,n})
                  \rtimes\mathfrak S_j\right).
\]
For \(b\ne0\), pullback identifies
\[
 H^j(X_{P_j},\kk)_b\simeq\bigwedge^j V_b,\qquad
 V_b=H^1(C,\kk)_b,
\]
with the exterior power embedded by antisymmetrization.
\end{lemma}
\begin{proof}
Each curve point has at most one zero coordinate.  Since \(j<m\),
the products are never all zero.  A product of evaluations of
linear polynomials is the evaluation of their product, so the
formula gives a morphism to \(X_{P_j}\).  Its defining sections
give
\[
 \mu_j^*\mathcal O(1)
 =\mathcal O_C(1)\boxtimes\cdots\boxtimes\mathcal O_C(1),
\]
which is ample.  A positive-dimensional projective fiber would
contain a curve on which this line bundle has degree zero.  Thus
the proper morphism is finite.

Homogenization \(f(z)\mapsto v^j f(u/v)\) identifies polynomials of
degree at most \(j\) with binary forms of degree \(j\).
On the bases the map multiplies
\(j\) binary linear forms.  Every degree-\(j\) form factors into such
forms, including the factor \(v\) representing a zero at infinity.
Choose roots of the factors' evaluation vectors; a phase change
on one factor then realizes any prescribed root vector above the
product.  This proves surjectivity.  Where the factors are
distinct and no evaluation vanishes, the only identifications
are their permutations and phase tuples of total product one.
The indicated finite quotient therefore maps finitely and
birationally to the smooth target.  Normality of the target
makes this map an isomorphism.

By Lemma~\ref{lem:quotient}, target cohomology is the invariant
cohomology of the product.  In residual character \(b\), product
kernel invariance forces character \(b\) in every factor.
Connectedness of \(C\) kills nontrivial characters in \(H^0,H^2\),
leaving \(V_b^{\otimes j}\).  A geometric interchange of two
degree-one factors is minus ordinary tensor interchange.
The permutation invariants are exactly the alternating tensors.
\end{proof}

\begin{proposition}[Character cohomology and support]\label{prop:characters}
For \(P\in\U_{d,m}\), invariant cohomology is the pullback of
\(H^*(\PP(P),\kk)\) and is spanned by algebraic hyperplane
classes.  Every nontrivial character occurs only in degree \(d\),
and
\[
 \dim_{\kk} W_{P,b}=\binom{s(b)-2}{d}\quad(b\ne0),
\]
where the value is zero if \(s(b)-2<d\).

If \(S=S(b)\) has \(s\geq d+2\) elements, projection identifies
\(P\) with a coordinate-general plane \(P_S\subset\C^S\)
and defines a finite quotient \(p_S:X_P\to X_{P_S}\).
Its pullback is an isomorphism
\(H^d(X_{P_S},\kk)_{b|_S}\simeq W_{P,b}\).
These descriptions commute with smooth parameter pullback.
\end{proposition}
\begin{proof}
The invariant assertion follows from the finite quotient formula.
The pullback of a base hyperplane is \(n\) times the root
hyperplane.  For \(d\geq1\), Lemma~\ref{lem:ambient} makes all
degrees outside \(d\) ambient and phase-invariant.  For \(d=0\),
degree zero is the regular representation of the finite torsor.

For \(d\geq1\), compute the rank at an evaluation plane.  The
curve \(C\to\PP^1\) is branched at the zeros of its \(m\)
coordinate forms.  Pass to the effective cyclic quotient for
\(b\).  Exactly the \(s=s(b)\) points in its support have
nontrivial inertia character.  Its character summand of the
finite pushforward of the constant sheaf is a rank-one local
system off those points, extended by zero at them.  Its
compactly supported Euler characteristic is \(2-s\), the
rank-one Euler characteristic of the punctured sphere.
The connected cyclic cover has no nontrivial character in
degrees zero or two.  Hence \(\dim V_b=s-2\).  The effective
quotient makes this valid for nonfaithful characters and
nonunit entries as well.

Lemma~\ref{lem:evaluation} gives the binomial rank there.  The
phase idempotent cuts out a local system on the connected
\(\U_{d,m}\), so the rank is constant.  In dimension zero every
character has rank one in the regular representation, agreeing
with the formula since a nonzero sum-zero tuple has \(s\geq2\).

For \(s\geq d+2\), any \(d+1\) selected coordinates detect \(P\).
Its projection is injective with coordinate-general image, and
the selected root coordinates cannot all vanish.  Forgetting
the others defines \(p_S\).  Its pullback of \(\mathcal O(1)\)
is \(\mathcal O_{X_P}(1)\), so the ample-line-bundle argument
proves finiteness.  It is surjective by taking roots of the
remaining coordinate values of the determined base point.

The fibers are exactly the orbits of
\(K_S=\ker(G_{m,n}\to G_{s,n})\).  For \(d\geq1\), the induced
finite map \(X_P/K_S\to X_{P_S}\) is birational and its target
is smooth, so it is an isomorphism.  For \(d=0\), it is directly
the quotient of finite torsors.  The character \(b\) is inflated
from \(G_{s,n}\), so Lemma~\ref{lem:quotient} proves the
eigenspace assertion.  Universal coordinate projections and
constant phase actions give the asserted family compatibility.
\end{proof}

\section{The perfect Gale correspondence}\label{sec:gale}
The rank formula makes the Fermat source of the Gale tensor a line.
We prove that its image pairs the two complementary character spaces
perfectly at every coordinate-general parameter.
The underlying orthogonal relation is linear Gale association;
compare \cite[preprint, Definition~1.1 and Proposition~2.2]{EP}.

\begin{theorem}[Perfect Gale tensor]\label{thm:gale}
Let \(P\in\U_{d,m}\), set
\[
 Q=P^\perp\quad\text{for }\sum_i u_i v_i,\qquad d'=m-d-2.
\]
Then \(Q\in\U_{d',m}\), and
\[
 f_P:X_P\times X_Q\longrightarrow F_n^{m-2},\qquad
 ([x_i],[y_i])\longmapsto[x_i y_i]
\]
is an algebraic morphism.  For every full-support character \(b\),
\(L_b=H^{m-2}(F_n^{m-2},\kk)_b\) is a line, and pullback maps it
to \(W_{P,b}\otimes W_{Q,b}\).  For every \(0\ne\ell\in L_b\), the
tensor \(f_P^*\ell\) has full pairing rank: the associated map
\(W_{P,b}^*\to W_{Q,b}\) is an isomorphism.  This includes \(d=0\)
and \(d'=0\).
\end{theorem}
\begin{proof}
We verify the morphism and its character, compute full rank at an
evaluation plane, transport that rank across the connected parameter
space, and then handle the zero-dimensional endpoints.

The maximal minors of the bilinear orthogonal complement are, up to
a common nonzero scalar and signs, the complementary maximal minors
of \(P\).  This follows by comparing top exterior vectors under the
standard volume pairing on \(\C^m\).  Hence \(Q\) is coordinate-general.

A point of \(X_P\) has at most \(d\) zero coordinates and a point
of \(X_Q\) at most \(d'\).  Their product vector has at most
\(m-2\) zeros, so it defines a morphism.  Its image is on the
Fermat hypersurface since
\(\sum_i(x_i y_i)^n=\sum_i x_i^n y_i^n=0\).
Equivariance for multiplication
\(G_{m,n}\times G_{m,n}\to G_{m,n}\) sends pullback character
\(b\) to \((b,b)\).  Proposition~\ref{prop:characters} leaves
only the two middle degrees.  Applied to the hyperplane
\(\sum_i z_i=0\), it also gives
\(\dim L_b=\binom{m-2}{m-2}=1\).

First suppose \(d,d'>0\), and put \(N=m-2=d+d'\).
Choose the evaluation planes of Lemma~\ref{lem:evaluation}, and set
\[
 c_i=\prod_{h\ne i}(a_i-a_h)^{-1}.
\]
Lagrange interpolation gives
\[
 \sum_i c_i a_i^e=0\qquad(0\leq e\leq m-2),
\]
by comparing the coefficient of \(z^{m-1}\) in the interpolation
formula for \(z^e\).  Products of polynomials of degrees at most
\(d,d'\) have degree at most \(m-2\); the identities and dimension
give
\[
 P_d^\perp=\operatorname{diag}(c_i)P_{d'},\qquad
 \{(z_i):\sum_i z_i=0\}=\operatorname{diag}(c_i)P_N.
\]
Choose \(\xi_i^n=c_i\), and let
\(s_\xi([z_i])=[\xi_i z_i]\) be coordinate scaling.
It restricts to phase-equivariant isomorphisms
\[
 s_\xi:X_{P_{d'}}\xrightarrow{\sim}X_{P_d^\perp},
 \qquad
 s_\xi:X_{P_N}\xrightarrow{\sim}F_n^N.
\]
The actual evaluation maps of Lemma~\ref{lem:evaluation} give
the commutative square
\[
\begin{array}{ccc}
 C^N=C^d\times C^{d'}
 &\xrightarrow{\ \mu_d\times(s_\xi\circ\mu_{d'})\ }&
 X_{P_d}\times X_{P_d^\perp}\\[3pt]
 {\scriptstyle s_\xi\circ\mu_N}\Big\downarrow
 &&\Big\downarrow{\scriptstyle f_{P_d}}\\[3pt]
 F_n^N&\xrightarrow{\ \id\ }&F_n^N .
\end{array}
\]
Indeed, both composites send an ordered tuple
\((x^{(1)},\ldots,x^{(N)})\) to
\([\xi_i\prod_{\nu=1}^N x_i^{(\nu)}]_{i=1}^m\).
With \(V=H^1(C,\kk)_b\), pullback in this square and
Lemma~\ref{lem:evaluation} identify the three character spaces as
\[
 L_b\simeq\bigwedge^N V,\qquad
 W_{P_d,b}\simeq\bigwedge^d V,\qquad
 W_{P_d^\perp,b}\simeq\bigwedge^{d'}V.
\]
Full support gives \(\dim V=N\).
Use the unnormalized alternating inclusion
\[
 \iota_j(v_1\wedge\cdots\wedge v_j)
 =\sum_{\sigma\in\mathfrak S_j}\operatorname{sgn}(\sigma)
 v_{\sigma(1)}\otimes\cdots\otimes v_{\sigma(j)}.
\]
More precisely, Lemma~\ref{lem:evaluation} says that \(\mu_j^*\)
identifies its character space with \(\operatorname{im}\iota_j\).
The exact equality of coordinate-product maps shows that the
induced pullback from the Fermat line is the exterior coproduct
\[
 \bigwedge^N V\longrightarrow\bigwedge^d V\otimes\bigwedge^{d'}V:
\]
after \(\iota_d\otimes\iota_{d'}\) it is \(\iota_N\), and these
inclusions are injective.  For a basis \(v_1,\ldots,v_N\), it is
\[
 v_1\wedge\cdots\wedge v_N\longmapsto
 \sum_{\substack{I\subset\{1,\ldots,N\}\\ |I|=d}}
 \operatorname{sgn}(I,I^c)\,v_I\otimes v_{I^c}.
\]
Each wedge uses increasing order.  Complementary wedge bases are
paired bijectively with nonzero coefficients, proving full rank
\(\binom{N}{d}\).

The map \(P\mapsto P^\perp\) is algebraic; the bilinear, not
Hermitian, orthogonal complement is essential here.  The maps
\(f_P\) form a morphism from the smooth projective product family
over \(\U_{d,m}\) to the fixed Fermat variety.  Cohomological
pullback is a morphism of local systems.  In flat character frames
the tensor is constant, so its matrix rank is locally constant.
Connectedness extends the computed rank to every \(P\).

If \(d=0\), choose a root point \(x\) above \(\PP(P)\).
Every \(x_i\ne0\), and multiplication by \(x\) identifies
\(X_Q\) with \(F_n^{m-2}\).  The restriction of \(f_P\) to
\(\{x\}\times X_Q\) is this isomorphism.  Its pullback on \(L_b\)
is nonzero.  A nonzero eigenfunction on the finite torsor \(X_P\)
is nonzero at every point, so evaluation at \(x\) detects a nonzero
\((b,b)\) tensor.  Both character ranks are one, hence it has
full rank.  The case \(d'=0\) is symmetric, including \(m=2\)
when all the varieties are finite.
\end{proof}

\subsection{A rational correspondence supplying cohomology}
The next objective is a forward source map to \(X_P\).  For a
support \(S\) with \(|S|\geq d+2\), project to \(P_S=\pr_S P\), then contract the
perfect Gale tensor against opposite-character cohomology of
\(Y_S=X_{P_S^\perp}\).  The source is \(F_n^{|S|-2}\times Y_S\).
The resulting graph cycle is rational, so the character calculation
after scalar extension will prove spanning for the cycle's action on
a whole rational source degree.

For a smooth projective source \(T\), a correspondence
\(\Gamma\in\CH^{\dim T+j}(T\times X)_\Q\) acts by
\[
 \Gamma_*\alpha=(\pr_X)_*(\pr_T^*\alpha\smile[\Gamma]).
\]
It sends \(H^{k-2j}(T,\Q)(-j)\) to \(H^k(X,\Q)\) as a Hodge
map.  Pullback by a morphism is represented by its transpose graph.

\begin{proposition}[Gale cohomology supply]\label{prop:supply}
Let \(P\in\U_{d,m}\).  For each \(S\subset\{1,\ldots,m\}\)
with \(s=|S|\geq d+2\), put
\[
 P_S=\pr_S P,\quad Q_S=P_S^\perp,\quad Y_S=X_{Q_S},
 \quad h_S=s-d-2.
\]
There is a rational algebraic correspondence from
\(F_n^{s-2}\times Y_S\) to \(X_P\), with degree shift
\(-2h_S\), whose image after scalar extension contains
\(W_{P,b}\) for every character with support \(S\).

Together with hyperplane cycles from a point, these finitely many
maps span the full rational cohomology of \(X_P\).  In the middle
degree the source maps may be taken on the entire groups
\[
 H^{d+2h_S}(F_n^{s-2}\times Y_S,\Q)(h_S).
\]
\end{proposition}
\begin{proof}
First target \(X=X_{P_S}\).  Put \(h=h_S\), \(F=F_n^{s-2}\),
and \(Y=Y_S\).  Thus \(\dim F=d+h\), \(\dim X=d\), and
\(\dim Y=h\).
The graph embedding
\[
 X\times Y\longrightarrow(F\times Y)\times X,\qquad
 (x,y)\longmapsto(f_{P_S}(x,y),y,x)
\]
defines \(\Gamma_S\in\CH^{d+h}(F\times Y\times X)_\Q\).
The source dimension is \(d+2h\), so the shift is \(-2h\).
On a K\"unneth tensor its action is
\[
 \alpha\otimes\beta\longmapsto
 (\pr_X)_*\bigl(f_{P_S}^*\alpha\smile\pr_Y^*\beta\bigr).
\]
For a character \(b\) with support \(S\), put \(b_S=b|_S\) and
\(L_{b_S}=H^{s-2}(F,\kk)_{b_S}\).
The invariant Poincar\'{e} pairing on \(Y\) pairs only opposite
characters and is perfect on each opposite pair, by its full
nondegeneracy and the direct character decomposition.  This also
holds for the sum pairing on a finite torsor.  It identifies
\(H^h(Y,\kk)_{-b_S}(h)\) with the dual of
\(H^h(Y,\kk)_{b_S}\).  Theorem~\ref{thm:gale} therefore makes the
displayed graph action an isomorphism of Hodge-graded spaces
\[
 L_{b_S}\otimes H^h(Y,\kk)_{-b_S}(h)
 \xrightarrow{\;\sim\;} H^d(X,\kk)_{b_S}.
\]
Composing with \(p_S^*\) fills \(W_{P,b}\).

Both the graph cycle and the transpose graph for \(p_S^*\) have
rational coefficients.  Apply them to the entire source degree,
without character projectors.  After scalar extension their images
contain all the character pieces just described.  Hyperplane cycles
from a point supply the invariant cohomology, including every degree
outside the middle.  The sum of these finite rational maps is thus
surjective after extension to \(\kk\); its rational cokernel is
already zero.  No algebraic K\"unneth projector is required.
\end{proof}

\begin{lemma}[Sources in codimension at most two]\label{lem:one-curve}
Let \(X_P\) have dimension \(d\) and codimension
\(c=m-d-1\in\{1,2\}\), and write \(S=S(b)\) for the support of a
character \(b\).
\begin{enumerate}[label=\textup{(\roman*)}]
\item If \(c=1\), every nontrivial character \(b\) with \(W_{P,b}\ne0\)
has full support, and its dual source \(Y_S\) in
Proposition~\ref{prop:supply} is finite.
\item If \(c=2\), every nontrivial character \(b\) with \(W_{P,b}\ne0\)
has support of size \(m\) or \(m-1\).  The full-support source is the one connected
curve \(C=X_{P^\perp}\).  If \(S=\{1,\ldots,m\}\setminus\{i\}\),
then \(Y_S\) is the reduced \(i\)th boundary fiber of \(C\), with
the zero coordinate omitted, and consists of \(n^{m-2}\) points.
\end{enumerate}
\end{lemma}
\begin{proof}
For \(b\ne0\), put \(S=S(b)\) and \(s=|S|\).  If \(W_{P,b}\ne0\),
then \(s\geq d+2\), and \(\dim Y_S=s-d-2\).  If \(c=1\), \(m=d+2\), so \(s=m\) is the
only possibility and the source dimension is zero.  If \(c=2\),
\(m=d+3\), so \(s=m\) gives dimension one and \(s=m-1\) gives
dimension zero.  The full-support source is \(X_{P^\perp}\),
connected by Lemma~\ref{lem:geometry}.

For a missing-one support, zero extension identifies
\[
 (\pr_S P)^\perp=\{q|_S:q\in P^\perp,\ q_i=0\}.
\]
This is the coordinate-general line cut from the dual two-plane
by its \(i\)th coordinate hyperplane.  Its root torsor is precisely
the reduced fiber \(\{y_i=0\}\) on \(C\) after forgetting that
coordinate.  The phase group has size \(n^{m-2}\).  Distinct
boundary fibers are disjoint because a root-curve point has at
most one zero coordinate.
\end{proof}

Finite dual sources require only point classes.  The only possible
contributing positive-dimensional dual source is therefore the
full-support curve in codimension two.  Its Fermat factor remains
in the cohomology supply.

\section{The full branch-configuration family}\label{sec:family}
For the curve of Lemma~\ref{lem:one-curve}, we compare every marked
cover fiber with transport in its normalized root family.  This will
verify that its tensor-generic parameters belong to
$\mathcal T_\phi^{\mathrm{Hg}}$, as required by
Theorem~\ref{p4:abelian-covers}.
The normalized diagonal curves are the standard generalized
Fermat models of \cite{GDL,HKLP}; the relative roots and based
marking needed here are specified below.

\subsection{The coefficient space as a product}
Let \(m\geq3\).  For a frame \(A=(a_1,\ldots,a_m)\) of
\(Q\in\D_{2,m}\), put \(\Delta_{ij}=\det(a_i,a_j)\).
The coordinate form on \(\PP(Q)\) is
\(\ell_i(s,t)=a_{1i}s+a_{2i}t\), with zero
\(p_i=[-a_{2i}:a_{1i}]\).  Let
\[
 M_m=\{(z_4,\ldots,z_m)\in(\A^1\setminus\{0,1\})^{m-3}:
 z_i\ne z_j\},
\]
with \(M_3\) a point.  Normalizing the first three labelled points
to \(\infty=[1:0]\), \(0=[0:1]\), \(1=[1:1]\) identifies
\(M_m\) with \(M_{0,m}\).

\begin{proposition}[The branch map]\label{prop:parameter}
The normalized branch configuration is the regular morphism
\[
 \lambda:\D_{2,m}\to M_m,\qquad
 z_i=\frac{\Delta_{2i}\Delta_{13}}{\Delta_{1i}\Delta_{23}}
 \quad(i\geq4).
\]
There is a global product isomorphism
\(\D_{2,m}\simeq M_m\times(\Gm)^{m-1}\) under which \(\lambda\)
is projection.  It is smooth and surjective, and the quotient by
the effective coefficient torus \((\Gm)^m/\Gm\) is \(M_m\).
\end{proposition}
\begin{proof}
The ratios are row-basis independent and their denominators are
units.  The Pl\"ucker relations give
\[
 z_i-1=\frac{\Delta_{12}\Delta_{3i}}{\Delta_{1i}\Delta_{23}},
 \qquad
 z_i-z_j=-\frac{\Delta_{12}\Delta_{13}\Delta_{ij}}
 {\Delta_{23}\Delta_{1i}\Delta_{1j}},
\]
so the values lie in \(M_m\).

Let \(A_0(z)\) have columns
\[
 (0,1)^t,\ (1,0)^t,\ (1,-1)^t,\
 (1,-z_4)^t,\ldots,(1,-z_m)^t.
\]
Set \(t_1=1\) and
\[
 t_2=\frac{\Delta_{23}}{\Delta_{13}},\quad
 t_3=\frac{\Delta_{23}}{\Delta_{12}},\quad
 t_i=\frac{\Delta_{1i}\Delta_{23}}{\Delta_{12}\Delta_{13}}
 \quad(i\geq4).
\]
All are nonzero regular functions of \(Q\).  With \(B=(a_1,a_2)\),
left multiplication by
\[
 H=\begin{pmatrix}0&t_2\\1&0\end{pmatrix}B^{-1}
\]
gives, by Cramer's rule,
\[
 Ha_1=(0,1)^t,\quad Ha_2=(t_2,0)^t,\quad
 Ha_3=(t_3,-t_3)^t,\quad Ha_i=(t_i,-t_i z_i)^t .
\]
Hence \(Q=\operatorname{row}(A_0(z)\operatorname{diag}(t_i))\).
Conversely, substituting this form into the ratios recovers
\(z,t_2,\ldots,t_m\).  These are inverse regular maps for the
product isomorphism.  The effective coefficient torus acts freely
and transitively on the second factor, giving the quotient claim.
\end{proof}

\begin{remark}
On framed coefficient matrices the map to \(\Conf_m(\PP^1)\)
is a principal \((\Gm)^m\)-bundle: fixing the zero points leaves
one independent nonzero scalar for each coordinate covector.
A row change \(A\mapsto gA\) acts on the points by \(g^{-t}\).
Thus the framed coefficient space surjects onto
\(\Conf_m(\PP^1)\).  After row changes are removed, the quotient
\(\D_{2,m}/((\Gm)^m/\Gm)\), not \(\D_{2,m}\) itself, is \(M_{0,m}\).
\end{remark}

\subsection{Coordinate roots and the fixed cover}
Put \(Q_0(z)=\operatorname{row}A_0(z)\) and
\(P_0(z)=Q_0(z)^\perp\).  The normal form gives
\[
 Q(t,z)=\operatorname{diag}(t_i)Q_0(z),\qquad
 P(t,z)=\operatorname{diag}(t_i^{-1})P_0(z).
\]
The normalized dual curves form a smooth projective family
\(\cY_0\to M_m\) by Lemma~\ref{lem:geometry}.
Let \(\iota_z:\PP^1\xrightarrow{\sim}\PP(Q_0(z))\) be the
projective isomorphism induced by the row frame \(A_0(z)\).
Write
\[
 q_z:=\iota_z^{-1}\circ q_{Q_0(z)}:X_{Q_0(z)}\longrightarrow\PP^1,
 \qquad D_z:=\{\infty,0,1,z_4,\ldots,z_m\}.
\]
The points of \(D_z\) retain this order and are regarded as a
reduced divisor.

\begin{proposition}[The normalized root family and its fixed deck datum]\label{prop:cover-family}
The map
\[
 \rho:\widetilde\D_{2,m,n}:=M_m\times(\Gm)^{m-1}\to\D_{2,m},
 \qquad (z,h_2,\ldots,h_m)\mapsto(z,h_2^n,\ldots,h_m^n)
\]
is a connected finite \'{e}tale cover of degree \(n^{m-1}\).
Set \(h_1=1\).  Over it, the algebraic maps
\([y_i]\mapsto[h_i y_i]\) and \([x_i]\mapsto[h_i^{-1}x_i]\)
identify the normalized dual and original root families with
the families for \(Q(t,z)\) and \(P(t,z)\), respectively.
They commute with \(G_{m,n}\).

Every \(q_z\) is a connected \(G_{m,n}\)-cover of \(\PP^1\),
with ordered branch points
\((\infty,0,1,z_4,\ldots,z_m)\).  Identify
\[
 G_{m,n}\simeq(\Z/n)^m/\langle(1,\ldots,1)\rangle
\]
using \(\zeta_n=\exp(2\pi i/n)\).  Its ordered positive local
monodromies are the images \((\bar e_1,\ldots,\bar e_m)\) of
the standard basis.  The family runs over all ordered branch
configurations with this fixed deck datum.
\end{proposition}
\begin{proof}
The map is the \(n\)-power map on a torus times the identity.
It is finite locally free of degree \(n^{m-1}\) and \'{e}tale
because \(n h_i^{n-1}\) is a unit.  Its source is irreducible.
The coordinate scalings satisfy the root equations by the
preceding formulas for \(P,Q\); their inverses and phase
equivariance are immediate.

On the normalized base \([s:t]\), the coordinate forms are
\[
 t,\ s,\ s-t,\ s-z_4t,\ldots,s-z_m t.
\]
Their ordered zeros are the stated branch points.  On \(t\ne0\),
put \(v=s/t\), \(w_i=y_i/y_1\).  The function field is
\[
 w_2^n=v,\qquad w_3^n=v-1,\qquad w_i^n=v-z_i\quad(i\geq4).
\]
Valuations at the distinct finite points show that these \(m-1\)
functions are independent in \(\C(v)^*/\C(v)^{*n}\).
Kummer theory gives degree \(n^{m-1}\).  Equivalently, adjoin
the roots successively and use an Eisenstein test at the new
branch point, where the earlier functions are units with local
roots.  Thus the cover is connected and its phase group is
its full deck group.

At a finite point indexed by \(i\geq2\), a positive meridian
multiplies \(w_i\) by \(\zeta_n\) and fixes the other local roots,
giving \(\bar e_i\).  At infinity all right sides have simple
poles, so a meridian multiplies every \(w_i\), \(i\geq2\), by
\(\zeta_n^{-1}\).  This is
\(\bar e_1=-\sum_{i\geq2}\bar e_i\).  The entries have order
\(n\), sum to zero, and generate the group.

For \(*\in\PP^1\setminus D_z\), positive meridians generate the
punctured-sphere fundamental group with their product as the
single relation.  The displayed
epimorphism is precisely abelianization modulo \(n\):
\[
 \pi_1(\PP^1\setminus D_z,*)\to H_1(\PP^1\setminus D_z,\Z/n)
 \simeq(\Z/n)^m/\langle(1,\ldots,1)\rangle .
\]
A pure orientation-preserving mapping class conjugates each positive
peripheral meridian, so it fixes this epimorphism.  The standard
classification of connected covers by epimorphisms shows that the
ordered tuple determines the single fixed cover type here.  For
each complex structure on the pointed base, pullback charts off
the punctures and the local ramification power charts give the
unique complex structure on its compact covering surface.  The
root equations realize it for every configuration in \(M_m\).
\end{proof}

The coefficient normalization and the coordinate-root cover have
now realized every normalized unmarked cover with its fixed deck
datum.  A based topological marking is additional data: it retains
an isotopy class relative to the labels.  We will reach each such
marking by a configuration path with a transported lift.
The following sphere-isotopy fact provides the required paths.

\begin{lemma}[Sphere isotopy]\label{lem:sphere-isotopy}
Every orientation-preserving self-homeomorphism of \(S^2\) is
isotopic to the identity through self-homeomorphisms, with no
marked points required to stay fixed during the isotopy.
\end{lemma}
\begin{proof}
We use the planar Schoenflies theorem in its annulus form:
the region between two nested Jordan curves is an annulus.
It gives the following disk extension fact.  If \(D_0,D_1\)
are closed disks in the interior of a closed disk \(U\), then
an orientation-preserving homeomorphism \(D_0\to D_1\) extends
to a homeomorphism of \(U\) which is the identity on \(\partial U\).
Indeed, the two complementary regions in \(U\) are annuli.
In annulus coordinates, the prescribed maps on their two boundary
circles extend across the annulus: orientation-preserving circle
homeomorphisms are isotopic, as is seen by taking increasing lifts
to \(\mathbb R\) and linearly interpolating the lifts.  Gluing this annulus
map to the given disk map proves the extension fact.

Let \(g\) be the given sphere homeomorphism.  Choose \(p\in S^2\),
a point \(q\) different from \(p\) and \(g(p)\), and a small closed disk \(D\)
about \(p\) for which both \(D\) and \(g(D)\) avoid \(q\).
Their compact union lies in the interior of a closed disk
\(U\subset S^2\setminus\{q\}\).  The disk extension fact gives
a homeomorphism \(f\) of \(U\), fixed on its boundary, which
agrees with \(g\) on \(D\).  Extend \(f\) by the identity outside
\(U\).  The Alexander trick isotopes \(f\) to the identity,
supported in \(U\): in unit-disk coordinates a boundary-fixed
map \(F\) is deformed by
\[
 F_t(x)=
 \begin{cases}
 tF(x/t),&|x|\leq t,\\
 x,&|x|\geq t,
 \end{cases}
 \qquad F_0=\id,\quad 0<t\leq1.
\]
Now \(f^{-1}g\) fixes \(D\) pointwise, so it is a boundary-fixed
homeomorphism of the complementary closed disk.  A second
Alexander trick isotopes it to the identity.  Composing the
two isotopies proves the assertion.
\end{proof}

To record a freely moving auxiliary basepoint, set
\[
 V_m=\{(z,b)\in M_m\times\PP^1:b\notin D_z\}\simeq M_{0,m+1}.
\]
Let \(B_{m,n}\) be the complement of the coordinate boundary in
the total dual curve family \(\cY_0\), and define
\(\beta:B_{m,n}\to V_m\) by \(\beta(z,y)=(z,q_z(y))\).
The two kinds of parameter maps are
\[
\begin{array}{ccccc}
 \widetilde\D_{2,m,n}
 &\xrightarrow[\text{forget roots}]{\rho}&
 \D_{2,m}
 &\xrightarrow[\text{branches}]{\lambda}&M_m\\[10pt]
 B_{m,n}
 &\xrightarrow[\text{forget lift}]{\beta}&
 V_m
 &\xrightarrow[\text{forget }b]{\pr_1}&M_m .
\end{array}
\]
The first row concerns coefficient normalization and finite
coordinate-root choices.  In the second, \(B_{m,n}\) records
a chosen point above \(b\).  A marking will be represented by a
path in \(V_m\) lifted to \(B_{m,n}\), not by a point of
\(\widetilde\D_{2,m,n}\).

\begin{lemma}[The based marked comparison]\label{lem:based}
The map \(\beta:B_{m,n}\to V_m\) defined above is a connected
finite \'{e}tale \(G_{m,n}\)-torsor.  Pullback of
the root family to \(B_{m,n}\) has the tautological chosen lift
\(y\) of the auxiliary basepoint \(b=q_z(y)\).

Use orientation-preserving homeomorphism markings, up to isotopy
relative to the labels.  Fix an initial based equivariant
identification with a root fiber.  For every marked complex
sphere with the datum of Proposition~\ref{prop:cover-family},
there is a path in \(V_m\), lifted from the chosen point of
\(B_{m,n}\), such that under its induced cohomological
identification the endpoint root Hodge filtration is exactly
the filtration of the prescribed marked cover.  Consequently,
a rational tensor which remains Hodge after root-family transport
along every path in \(M_m\) is Hodge in every marked fiber.

The parameter space making both coordinate-root and based-lift
choices is
\[
 \widetilde\D_{2,m,n}\times_{M_m}B_{m,n}
 \simeq B_{m,n}\times(\Gm)^{m-1}.
\]
It is connected and projects smoothly and surjectively to
\(B_{m,n}\).
\end{lemma}
\begin{proof}
Away from \(D_z\), \(q_z\) is finite \'{e}tale with
free transitive phase fibers.  This proves the torsor assertion
except for connectedness.  The total \(\cY_0\) is irreducible:
it is smooth over the irreducible \(M_m\), and every fiber is
connected.  Indeed its irreducible components, being disjoint
open components of a smooth variety, would have nonempty open
images on the base; two such images meet, contradicting
connectedness of that fiber.  Removing the boundary leaves a
nonempty irreducible open.  Thus \(B_{m,n}\) is connected.

Fix \((z_0,y_0)\in B_{m,n}\), put \(b_0=q_{z_0}(y_0)\), and
identify the fixed labelled topological sphere with the normalized
base at \((z_0,b_0)\).  Also fix a \(G_{m,n}\)-equivariant
identification of its topological cover with the root cover,
taking the fixed upstairs point to \(y_0\).  The based
epimorphisms agree because both take the ordered positive meridians
to \(\bar e_i\).  A change of the chosen lift conjugates the
epimorphism and hence leaves it unchanged for this abelian group.

Consider any prescribed marked complex sphere.  A complex curve
of genus zero is \(\PP^1\), and its first three branch labels
give a unique M\"obius normalization to \(\infty,0,1\).
The marking is then represented by an orientation-preserving
self-homeomorphism \(g\) of the reference sphere which carries
all reference labels to the normalized target labels.  It fixes
the first three labels at the endpoint.  By
Lemma~\ref{lem:sphere-isotopy}, choose an unmarked isotopy \(g_t\)
from the identity to \(g\).  For each \(t\), let \(a_t\) be the
unique M\"obius transformation taking
\(g_t(p_1),g_t(p_2),g_t(p_3)\) to \(\infty,0,1\).
The formulas for the transformation of a distinct triple show
that \(a_t\) is continuous, and \(a_0=a_1=\id\).
Thus \(a_tg_t\) keeps the first three images fixed.  Its remaining
branch images and its auxiliary-point image give a path
\(\gamma\) in \(V_m\), starting at \((z_0,b_0)\), whose endpoint
marking is the prescribed \(g\), not only a marking of the same
unmarked fiber.

Lift \(\gamma\) uniquely to \(B_{m,n}\) from \((z_0,y_0)\);
write \(y_t\) for its tautological upstairs point.  Naturality
of oriented labelled meridians under \(a_tg_t\) and the
abelianization description in Proposition~\ref{prop:cover-family}
give equality of the exact based epimorphisms at every \(t\).
The covering-space lifting criterion therefore gives a unique
\(G_{m,n}\)-equivariant lift of the base marking which sends
the fixed upstairs point to \(y_t\).  Covering homotopy lifting
and uniqueness make these lifts continuous along \(\gamma\).
They extend over the compactification through the fixed local
ramification power charts, giving a continuous trivialization of
the compact root family pulled back to the path interval.
They therefore induce exactly its cohomological parallel transport
along the projection of \(\gamma\) to \(M_m\).

At the endpoint, this lifted covering identification is
holomorphic for the complex structure of the prescribed marked
cover: off the branch points both complex structures are pulled
back from the normalized base, and the same statement holds in
the local power charts.  Hence it identifies the endpoint Hodge
filtrations.  This proves the asserted tensor consequence.
No loop in \(V_m\) is required to lift to a loop in \(B_{m,n}\);
its endpoint lift may differ by a deck phase, which is precisely
retained by the transported point \(y_t\).

Finally \(V_m\to M_m\) is smooth and surjective, with fiber the
punctured line.  Its composite with the finite \'{e}tale torsor
is smooth and surjective.  The product description of
\(\widetilde\D_{2,m,n}\) gives the displayed fiber product;
it is connected and its projection is the torus projection.
\end{proof}

\begin{remark}[Coordinate roots and automorphisms]
The cover \(\rho\) records coordinate-root choices.  It cannot
generally be omitted from a global algebraic family isomorphism:
a torus coordinate has no regular \(n\)th root on the original
torus.  Different choices change the dual isomorphism by a
deck phase.  More precisely, let \(T_{\rm root}=(\Gm)^m/\Gm\)
act on root coordinates and let \(T_{\rm coef}=(\Gm)^m/\Gm\)
be the effective coefficient torus of Proposition~\ref{prop:parameter}.
For the dual curve the induced map is
\([n]:T_{\rm root}\to T_{\rm coef}\); the inverse coordinate
convention for the original intersection composes it with inversion.
Its kernel is \(G_{m,n}\).  Thus \(T_{\rm coef}\) acts freely
on the coefficient locus, whereas the \(T_{\rm root}\)-action
retains exactly the phase stabilizer.

\end{remark}

\subsection{Algebraic very-generality}
The tensor condition in the marked-cover theorem is phrased on a
marked parameter space.
For the present algebraic family its exceptional tensor loci are
algebraic on \(M_m\), which gives the algebraic meaning of
\emph{very general} in Theorem~\ref{thm:main}.

\begin{lemma}[Tensor-generic points]\label{lem:tensor-generic}
Let \(S\) be a smooth connected complex algebraic variety carrying
an integral polarized variation \(\V\) of weight one.  On the
universal cover of \(S^{\rm an}\), mark its underlying local system.
Outside a countable union of proper closed algebraic subsets of \(S\),
every rational tensor in \(\V_s^{\otimes 2a}\) of type \((a,a)\)
remains of that type after parallel transport along every path
starting at \(s\), for every \(a\geq0\).
\end{lemma}
\begin{proof}
The tensor variation \(\V^{\otimes 2a}(a)\), with the standard
tensor/Tate polarization rescaled by a positive integer if needed,
is integral, polarized, and of weight zero.  Denote its polarization
by \(Q_a\).
Multiplying a denominator makes a rational tensor integral.
For an integral flat tensor \(u\) on the universal cover, let
\(H_u\subset S\) be the set where some determination of \(u\)
has type \((0,0)\).  Cattani--Deligne--Kaplan,
Corollary~1.3 on p.~2 of the preprint \cite{CDK}, says
that \(H_u\) is algebraic.  Its proof writes \(H_u\), for
\(K=Q_a(u,u)\), as a union of images of connected components of
the algebraic variety \(S^{(K)}\) finite over \(S\) in their
Theorem~1.1.  There are finitely many such components, and finite
morphisms are closed, so \(H_u\) is closed.

If \(u\) is not Hodge throughout the universal cover, \(H_u\)
is proper.  On a simply connected analytic chart, its determinations
form a countable set of integral tensors.  Each has a closed analytic
Hodge locus: for a real weight-zero tensor, type \((0,0)\) is
equivalent to membership in the holomorphic subbundle \(F^0\).
None contains an open subset, since analytic continuation on the
connected universal cover and deck equivariance would then make
\(u\) globally Hodge.  If \(H_u=S\), the chart would be covered
by countably many nowhere-dense closed analytic subsets, contrary
to Baire's theorem.  Hence \(H_u\) is proper.

There are countably many integral tensors in all powers.  Remove
\(H_u\) for the ones which are not globally Hodge.  A rational
Hodge tensor at a remaining point, after clearing a denominator,
must be globally Hodge on the universal cover, equivalently
Hodge after parallel transport along every path from that point.
\end{proof}

\begin{corollary}\label{cor:good-curves}
There is a countable union \(Z\) of proper closed algebraic subsets
of \(M_m\) such that
for \(z\notin Z\), every CM abelian \(B\), and every \(N\geq1\),
the rational Hodge conjecture holds on
\((\Jac(X_{Q_0(z)})\times B)^N\).  The same holds for the dual
curve at every point of \(\D_{2,m}\setminus\lambda^{-1}(Z)\).
\end{corollary}
\begin{proof}
Write \(\pi:\cY_0\to M_m\).  Its integral polarized variation is
\(\V=R^1\pi_*\Z\).  Apply Lemma~\ref{lem:tensor-generic}.
For a remaining \(z\), choose an auxiliary point and one of its
lifts, hence a point of \(B_{m,n}\), and fix the initial based
identification.  A rational tensor Hodge in this marked fiber
corresponds to a rational tensor in the root fiber.  By the
lemma, it remains Hodge under transport along every path in
\(M_m\).  For any prescribed marked fiber,
Lemma~\ref{lem:based} supplies a path in \(V_m\) and its based
lift whose endpoint filtration is that of the prescribed cover.
Applying root-family globality to the projected path shows that
the tensor is Hodge in that marked fiber.  The fiber was arbitrary,
so this marked parameter belongs to $\mathcal T_\phi^{\mathrm{Hg}}$.
Theorem~\ref{p4:abelian-covers} applies.

Surjectivity of \(\lambda\) makes inverse images of the proper
algebraic exceptional subsets proper.  At a complex point above
\(z\), roots of its torus coordinates exist, and Proposition~
\ref{prop:cover-family} identifies its dual curve with
\(X_{Q_0(z)}\).  The Hodge conclusion is invariant under this
algebraic isomorphism.
\end{proof}

\section{Hodge classes on all self-powers}\label{sec:assembly}
Lemma~\ref{lem:one-curve} leaves at most one contributing
positive-dimensional dual source.  Section~\ref{sec:family}
supplies its cover-Jacobian conclusion; we now combine it with
CM sources for the Fermat factors.

\begin{proof}[Proof of Theorem~\ref{thm:main}]
For \(c=0\), products of hyperplane classes span every power of
projective space.  If \(d=0\), the smooth target and all its powers
are reduced finite sets, whose point cycles span cohomology.
Assume henceforth \(d\geq1\) and \(c=1\) or \(2\).

Proposition~\ref{prop:supply} gives finitely many rational algebraic
maps spanning each cohomology degree of \(X_P\).  For a nontrivial
character with support \(S\), their source is
\(F_n^{s-2}\times Y_S\), with shift \(-2h_S\).  Apply
Proposition~\ref{p4:fermat-supply} to
\(H^{s-2}(F_n^{s-2},\Q)\), the cohomology group that witnesses
character spanning.  Compose its finite
CM-source maps with the Gale correspondence.  These compositions
are rational algebraic: pull the cycles to a triple product,
intersect them, and push forward.

For a finite \(Y_S\), use the inclusions of all its points to
span its entire degree-zero cohomology.  For the full-support
curve \(C\), choose an Abel map \(C\to\Jac(C)\).  Its pullback
on \(H^1\) is an isomorphism induced by the transpose graph,
and supplies the factor required in the contraction.  Hyperplane
cycles from a point supply the ambient target classes.

Suppose one Fermat-source map from a CM abelian variety \(B\)
has shift \(2j\), and put \(h=h_S\).
Let \(A_S=\Jac(C)\) when \(Y_S=C\), and let \(A_S\) be a point
for one chosen component when \(Y_S\) is finite.  The composite with the
degree-preserving \(A_S\)-to-\(Y_S\) map and the Gale shift
\(-2h\) has total shift \(2(j-h)\):
\[
 H^{d+2h-2j}(B\times A_S,\Q)(h-j)
       \longrightarrow H^d(X_P,\Q).
\]
Its K\"unneth subspace
\(H^{d+h-2j}(B)\otimes H^h(A_S)\) is the part used to witness
spanning.  The map itself is applied to the whole source degree.

Put \(A_C=\Jac(C)\) in codimension two and let \(A_C\) be a point
in codimension one.  Combining these maps with the ambient point
correspondences gives, in every degree \(k\), a surjection
\[
 \bigoplus_\nu
 H^{k-2j_\nu}\bigl(A_C^{v_\nu}\times B_\nu,\Q\bigr)(-j_\nu)
 \longrightarrow H^k(X_P,\Q),                         \tag{*}
\]
where \(v_\nu\in\{0,1\}\) and \(B_\nu\) is CM.  In codimension
one all \(v_\nu=0\).  The integers \(j_\nu\) include the negative
Gale shifts and the CM-source shifts.  All maps are rational
algebraic correspondences.  After scalar extension to \(\kk\), the
specified K\"unneth character pieces span by Proposition~\ref{prop:supply}.
Applying each rational correspondence to the whole source degree
only enlarges its image, so the rational
cokernel vanishes.  No algebraic K\"unneth projector is used.

It remains to verify the finite-product Hodge hypothesis for these
sources; Companion results~\ref{comp:transfer} will then
apply to \((*)\) and give all self-powers of \(X_P\) at once.
For \(c=1\), take as source collection all CM abelian varieties,
including a point.  Their finite products are CM, so the CM Hodge
theorem \cite[Theorem~1.1]{OpenAICM} gives the required hypothesis
at every parameter.  Equivalently, coordinate scaling identifies
each smooth one-equation diagonal hypersurface with a Fermat hypersurface.

For \(c=2\), choose \(P^\perp\) outside the union in
Corollary~\ref{cor:good-curves}.  For this fixed curve \(C\), let
\[
 \mathscr S_C=\{\Jac(C)^a\times M:a\geq0,\ M\text{ a CM abelian variety}\}.
\]
A finite product of these sources again has the displayed form.
If \(a\geq1\), it is an algebraic retract of
\((\Jac(C)\times M)^a\): retain every Jacobian factor and one
\(M\) factor, and set the other \(M\) coordinates to zero for the
section.  Pull a Hodge class on the retract to that power.
Corollary~\ref{cor:good-curves} makes it algebraic there, and pullback
along the section returns an algebraic representative on the retract.
The corollary applies to every \(M\) at the same parameter.  If \(a=0\),
the CM Hodge theorem applies.
Thus \(\mathscr S_C\) satisfies the full finite-product hypothesis of
Companion results~\ref{comp:transfer} as well.  Applying that
transfer to \((*)\), with target \(X_P\), proves the assertion for
every self-power in both cases.

The exceptional union depends only on the tensor variation of
the single curve, not on the number of factors or the CM sources.
The same parameter therefore works for every self-power.  Proposition~
\ref{prop:smooth-locus} identifies its ambient open with the
full smooth diagonal parameter locus.
\end{proof}

\begin{remark}[Low-dimensional checks]
The dual curve has degree \(n^{m-1}\) over the line and \(m\)
branch points of inertia order \(n\).  Riemann--Hurwitz gives
\[
 2g(C)-2=n^{m-2}\bigl(m(n-1)-2n\bigr).
\]
For \(m=3\), the original codimension-two target is \(n^2\)
points and the dual is a three-branch Fermat curve of genus
\((n-1)(n-2)/2\); \(M_3\) is a point.  For \(m=4,n=2\), the
dual has genus one and still has the full four-point cover
variation.  If \(n=2\) and \(m\) is odd, the only candidate
full-support tuple has odd coordinate sum, so only finite
sources occur.
\end{remark}

\begin{thebibliography}{99}
\addcontentsline{toc}{section}{\refname}

\bibitem{Andre1992FixedPart}
Y.~Andr\'e,
\href{https://www.numdam.org/article/CM_1992__82_1_1_0.pdf}{\emph{Mumford--Tate groups of mixed Hodge structures and the theorem of the fixed part}},
Compositio Mathematica \textbf{82} (1992), no.~1, 1--24.

\bibitem{AokiCI}
N.~Aoki,
\emph{A note on complete intersections of Fermat type},
Commentarii Mathematici Universitatis Sancti Pauli
\textbf{35} (1986), 231--245.
\href{https://rikkyo.repo.nii.ac.jp/records/10098}{Institutional record}.

\bibitem{AokiFermat}
N.~Aoki,
\emph{Some new algebraic cycles on Fermat varieties},
Journal of the Mathematical Society of Japan \textbf{39} (1987),
385--396.
\href{https://doi.org/10.2969/jmsj/03930385}{doi:10.2969/jmsj/03930385}.

\bibitem{CDK}
E.~Cattani, P.~Deligne, and A.~Kaplan,
\emph{On the locus of Hodge classes},
preprint arXiv:alg-geom/9402009v1, 10 February 1994,
\href{https://publications.ias.edu/sites/default/files/70_LocusHodgeClasses.pdf}{preprint PDF}.
Published in
Journal of the American Mathematical Society \textbf{8} (1995),
483--506.
\href{https://doi.org/10.1090/S0894-0347-1995-1273413-2}{doi:10.1090/S0894-0347-1995-1273413-2}.

\bibitem{Deligne1974}
P.~Deligne,
\href{https://doi.org/10.1007/BF02685881}{\emph{Th\'eorie de Hodge, III}},
Publications Math\'ematiques de l'IH\'ES \textbf{44} (1974), 5--77.

\bibitem{DeligneMostow1986}
P.~Deligne and G.~D.~Mostow,
\href{https://doi.org/10.1007/BF02831622}{\emph{Monodromy of hypergeometric functions and non-lattice integral monodromy}},
Publications Math\'ematiques de l'IH\'ES \textbf{63} (1986), 5--89.

\bibitem{DO}
I.~Dolgachev and D.~Ortland,
\emph{Point sets in projective spaces and theta functions},
Ast\'erisque \textbf{165}, Soci\'et\'e math\'ematique de France, 1988.
\href{https://www.numdam.org/item/AST_1988__165__1_0/}{Numdam record}.

\bibitem{EP}
D.~Eisenbud and S.~Popescu,
\emph{The projective geometry of the Gale transform},
Journal of Algebra \textbf{230} (2000), 127--173.
\href{https://doi.org/10.1006/jabr.1999.7940}{doi:10.1006/jabr.1999.7940};
\href{https://arxiv.org/pdf/math/9807127}{preprint arXiv:math/9807127}.

\bibitem{Etingof2024Lie}
P.~Etingof,
\emph{Lie groups and Lie algebras},
\href{https://arxiv.org/abs/2201.09397v4}{arXiv:2201.09397v4} (2024).
Section locators refer to this version.

\bibitem{FarbMargalit2012}
B.~Farb and D.~Margalit,
\href{https://doi.org/10.1515/9781400839049}{\emph{A Primer on Mapping Class Groups}},
Princeton Mathematical Series 49, Princeton University Press, 2012.

\bibitem{GDL}
G.~Gonz\'alez-Diez, R.~A.~Hidalgo, and M.~Leyton,
\emph{Generalized Fermat curves},
Journal of Algebra \textbf{321} (2009), 1643--1660.
\href{https://doi.org/10.1016/j.jalgebra.2009.01.002}{doi:10.1016/j.jalgebra.2009.01.002};
\href{https://verso.mat.uam.es/~gabino.gonzalez/GenFerCurves.pdf}{author proof}.

\bibitem{Griffiths1968PeriodsII}
P.~A.~Griffiths,
\href{https://doi.org/10.2307/2373485}{\emph{Periods of integrals on algebraic manifolds, II: Local study of the period mapping}},
American Journal of Mathematics \textbf{90} (1968), no.~3, 805--865.

\bibitem{HKLP}
R.~A.~Hidalgo, A.~Kontogeorgis, M.~Leyton-\'Alvarez, and
P.~Paramantzoglou,
\emph{Automorphisms of generalized Fermat curves},
Journal of Pure and Applied Algebra \textbf{221} (2017), 2312--2337.
\href{https://doi.org/10.1016/j.jpaa.2016.12.011}{doi:10.1016/j.jpaa.2016.12.011};
\href{https://arxiv.org/pdf/1409.3063}{preprint arXiv:1409.3063}.

\bibitem{LandesmanLittSawin2025}
A.~Landesman, D.~Litt, and W.~Sawin,
\href{https://doi.org/10.2140/gt.2025.29.2733}{\emph{Big monodromy for higher Prym representations}},
Geometry \& Topology \textbf{29} (2025), no.~5, 2733--2782.
Theorem locators refer to
\href{https://arxiv.org/abs/2401.13906v2}{arXiv:2401.13906v2}.

\bibitem{Milnor}
J.~Milnor,
\emph{Morse Theory},
Annals of Mathematics Studies 51,
Princeton University Press, 1963.
\href{https://math.stanford.edu/~ralph/math215b/Milnor.pdf}{Public scan}.

\bibitem{Moonen2010}
B.~Moonen,
\emph{Special subvarieties arising from families of cyclic covers of
the projective line},
Documenta Mathematica \textbf{15} (2010), 793--819.
\href{https://doi.org/10.4171/DM/314}{doi:10.4171/DM/314}.
Locators refer to
\href{https://arxiv.org/pdf/1006.3400v2}{arXiv:1006.3400v2} (2010).

\bibitem{OpenAICM}
OpenAI,
\emph{The rational Hodge conjecture for CM abelian varieties},
OpenAI Math Release preprint
\href{https://github.com/openai/math/blob/main/preprints/The-rational-Hodge-conjecture-for-CM-abelian-varieties-September-30-2026/paper.pdf}{OAI:The-rational-Hodge-conjecture-for-CM-abelian-varieties-September-30-2026},
2026.

\bibitem{OpenAIAbelianPowers}
OpenAI,
\emph{Weil classes and Hodge classes on abelian powers},
OpenAI Math Release preprint
\href{https://github.com/openai/math/blob/main/preprints/Weil-classes-and-Hodge-classes-on-abelian-powers-September-30-2026/paper.pdf}{OAI:Weil-classes-and-Hodge-classes-on-abelian-powers-September-30-2026},
2026.

\bibitem{Rohde2009}
J.~C.~Rohde,
\emph{Cyclic coverings, Calabi--Yau manifolds and complex multiplication},
Lecture Notes in Mathematics \textbf{1975}, Springer, Berlin, 2009.
The preprint locator refers to
\href{https://arxiv.org/pdf/0711.2195v4}{arXiv:0711.2195v4} (2008).

\bibitem{Schoen}
C.~Schoen,
\emph{Hodge classes on self-products of a variety with an automorphism},
Compositio Mathematica \textbf{65} (1988), 3--32.
\href{https://www.numdam.org/article/CM_1988__65_1_3_0.pdf}{Original article}.

\bibitem{SK}
T.~Shioda and T.~Katsura,
\emph{On Fermat varieties},
Tohoku Mathematical Journal, Second Series \textbf{31} (1979),
97--115.
\href{https://doi.org/10.2748/tmj/1178229881}{doi:10.2748/tmj/1178229881}.

\bibitem{Terasoma}
T.~Terasoma,
\emph{Complete intersections of hypersurfaces---the Fermat case
and the quadric case},
Japanese Journal of Mathematics, New Series \textbf{14} (1988),
309--384.
\href{https://doi.org/10.4099/math1924.14.309}{doi:10.4099/math1924.14.309}.

\end{thebibliography}
\end{document}
